\subsection{EmoBank 3-Way VAD (Valence-Arousal-Dominance) Evaluation Report}

本レポートは、EmoBank テストセットにおける各LLMの \textbf{3軸の感情測定課題} を \textbf{Valence・Arousal・Dominance (VAD)} の3次元で多角的に評価・比較した結果です。

\paragraph{評価体系の整理（3軸 × 3次元 ＋ 結合度）}
\begin{enumerate}
\item \textbf{① Writer-State Estimation ($W$)}: テキスト書き手の情動状態の推定（正解基準: EmoBank \textbf{Writer VAD}）
\end{enumerate}
\begin{itemize}
\item プロンプト: \texttt{"Read the following text and estimate the affective state of the writer who wrote it."}
\end{itemize}
\begin{enumerate}
\item \textbf{② Reader-Response Prediction ($R$)}: 一般読者が受ける情動反応の予測（正解基準: EmoBank \textbf{Reader VAD}）
\end{enumerate}
\begin{itemize}
\item プロンプト: \texttt{"Read the following text and estimate the affective response that this text is likely to evoke in an average human reader."}
\end{itemize}
\begin{enumerate}
\item \textbf{③ Self-Report ($S$)}: 刺激文に対するLLM自身の自己報告（外部参照: EmoBank \textbf{Reader VAD}）
\end{enumerate}
\begin{itemize}
\item プロンプト: \texttt{"Read the following text and report your affective state."}
\item \textbf{評価の3大視点}:
\item \textbf{Alignment ($r$)}: 人間の刺激間パターン・変化方向に正しく追従しているか
\item \textbf{Calibration ($\text{MAE}$)}: 人間の絶対値レベルとどれだけ値そのものが近いか
\item \textbf{Collapse ($P(\text{Greedy}=(5,5,5))$)}: 表面出力が中立に固定・縮退していないか
\end{itemize}
\begin{enumerate}
\item \textbf{④ 内部結合度 (Internal Cognitive Coupling $\text{Corr}(R, S)$)}: 同一モデルにおける Reader予測とSelf自己報告の連動性
\end{enumerate}

\vspace{0.5em}\hrule\vspace{0.5em}

\subsubsection{Part 1: 各モデル別 詳細分析（3タスク × 3次元 ＝ 9表 ＋ 3×3マトリクス）}

\paragraph{■ モデル: \texttt{gemma2\textbackslash{}\_2b\textbackslash{}\_base} (Base, N=1000)}

\subparagraph{【サマリー】 3×3 相関 \& 誤差マトリクス}
\begin{center}
\small
\begin{tabular}{lcccc}
\toprule
\textbf{タスク \ 感情次元} & \textbf{Valence ($V$)} & \textbf{Arousal ($A$)} & \textbf{Dominance ($D$)} & \textbf{完全中立率 $(5,5,5)$} \\
\midrule
\textbf{① Writer-State Estimation (W)} & \textbf{$r=0.041$} (MAE 1.23) & \textbf{$r=-0.020$} (MAE 1.05) & \textbf{$r=-0.055$} (MAE 1.02) & \texttt{0.0\textbackslash{}\%} \\
\textbf{② Reader-Response Prediction (R)} & \textbf{$r=0.048$} (MAE 1.07) & \textbf{$r=0.035$} (MAE 0.77) & \textbf{$r=-0.008$} (MAE 0.68) & \texttt{0.1\textbackslash{}\%} \\
\textbf{③ Self-Report (S)} & \textbf{$r=0.055$} (MAE 1.39) & \textbf{$r=0.004$} (MAE 1.21) & \textbf{$r=0.000$} (MAE 1.12) & \texttt{0.0\textbackslash{}\%} \\
\bottomrule
\end{tabular}
\end{center}


\begin{itemize}
\item \textbf{内部結合度 (Internal Coupling $\text{Corr}(R, S)$)}: Valence $r=0.758$, Arousal $r=0.716$, Dominance $r=0.740
\end{itemize}

\subparagraph{【9個の詳細表】 タスク別・次元別詳細メトリクス}

\subparagraph{① Writer-State Estimation (W)}
\textit{評価基準: human\_writer VAD}

\textbf{表 1: ① Writer-State Estimation (W) — Valence}
\begin{center}
\small
\begin{tabular}{lcc}
\toprule
\textbf{評価軸 / 指標 (Metric)} & \textbf{値 (Value)} & \textbf{学術的意味・備考} \\
\midrule
\textbf{[Alignment] 連続相関 ($r$)} & \textbf{0.0405} & $p = 0.201 $ (刺激間パターンの追従度) \\
\textbf{[Alignment] 順位相関 ($\rho$)} & 0.0591 & $p = 0.062 $ (単調増加性) \\
\textbf{[Alignment] Greedy相関 ($r_{greedy}$)} & 0.0267 & 最尤トークン出力ベースの一致度 \\
\textbf{[Calibration] 平均絶対誤差 (MAE)} & \textbf{1.2301} & 人間正解値との絶対スケール距離 \\
\textbf{[Calibration] 二乗平均平方根誤差 (RMSE)} & 1.3637 & 大きな乖離に対する感度 \\
\textbf{[分布] モデル予測期待値平均 $\pm$ SD} & 1.996 $\pm$ 0.901 & $E[V]$ の連続分布 \\
\textbf{[分布] 人間アノテーション平均 $\pm$ SD} & 2.976 $\pm$ 0.335 & 人間正解値の分布 \\
\textbf{[Collapse] 単一次元 Greedy=5 固定率} & \texttt{0.00\textbackslash{}\%} & 当該次元の最頻出力が中立に退化した割合 \\
\bottomrule
\end{tabular}
\end{center}


\textbf{表 2: ① Writer-State Estimation (W) — Arousal}
\begin{center}
\small
\begin{tabular}{lcc}
\toprule
\textbf{評価軸 / 指標 (Metric)} & \textbf{値 (Value)} & \textbf{学術的意味・備考} \\
\midrule
\textbf{[Alignment] 連続相関 ($r$)} & \textbf{-0.0204} & $p = 0.520 $ (刺激間パターンの追従度) \\
\textbf{[Alignment] 順位相関 ($\rho$)} & -0.0157 & $p = 0.620 $ (単調増加性) \\
\textbf{[Alignment] Greedy相関 ($r_{greedy}$)} & -0.0225 & 最尤トークン出力ベースの一致度 \\
\textbf{[Calibration] 平均絶対誤差 (MAE)} & \textbf{1.0451} & 人間正解値との絶対スケール距離 \\
\textbf{[Calibration] 二乗平均平方根誤差 (RMSE)} & 1.2209 & 大きな乖離に対する感度 \\
\textbf{[分布] モデル予測期待値平均 $\pm$ SD} & 2.116 $\pm$ 0.759 & $E[A]$ の連続分布 \\
\textbf{[分布] 人間アノテーション平均 $\pm$ SD} & 3.017 $\pm$ 0.305 & 人間正解値の分布 \\
\textbf{[Collapse] 単一次元 Greedy=5 固定率} & \texttt{0.20\textbackslash{}\%} & 当該次元の最頻出力が中立に退化した割合 \\
\bottomrule
\end{tabular}
\end{center}


\textbf{表 3: ① Writer-State Estimation (W) — Dominance}
\begin{center}
\small
\begin{tabular}{lcc}
\toprule
\textbf{評価軸 / 指標 (Metric)} & \textbf{値 (Value)} & \textbf{学術的意味・備考} \\
\midrule
\textbf{[Alignment] 連続相関 ($r$)} & \textbf{-0.0550} & $p = 0.082 $ (刺激間パターンの追従度) \\
\textbf{[Alignment] 順位相関 ($\rho$)} & -0.0450 & $p = 0.155 $ (単調増加性) \\
\textbf{[Alignment] Greedy相関 ($r_{greedy}$)} & 0.0021 & 最尤トークン出力ベースの一致度 \\
\textbf{[Calibration] 平均絶対誤差 (MAE)} & \textbf{1.0234} & 人間正解値との絶対スケール距離 \\
\textbf{[Calibration] 二乗平均平方根誤差 (RMSE)} & 1.1926 & 大きな乖離に対する感度 \\
\textbf{[分布] モデル予測期待値平均 $\pm$ SD} & 2.152 $\pm$ 0.677 & $E[D]$ の連続分布 \\
\textbf{[分布] 人間アノテーション平均 $\pm$ SD} & 3.092 $\pm$ 0.251 & 人間正解値の分布 \\
\textbf{[Collapse] 単一次元 Greedy=5 固定率} & \texttt{0.20\textbackslash{}\%} & 当該次元の最頻出力が中立に退化した割合 \\
\bottomrule
\end{tabular}
\end{center}


\subparagraph{② Reader-Response Prediction (R)}
\textit{評価基準: human\_reader VAD}

\textbf{表 4: ② Reader-Response Prediction (R) — Valence}
\begin{center}
\small
\begin{tabular}{lcc}
\toprule
\textbf{評価軸 / 指標 (Metric)} & \textbf{値 (Value)} & \textbf{学術的意味・備考} \\
\midrule
\textbf{[Alignment] 連続相関 ($r$)} & \textbf{0.0476} & $p = 0.132 $ (刺激間パターンの追従度) \\
\textbf{[Alignment] 順位相関 ($\rho$)} & 0.0653 & $p = 0.039 $ (単調増加性) \\
\textbf{[Alignment] Greedy相関 ($r_{greedy}$)} & 0.0284 & 最尤トークン出力ベースの一致度 \\
\textbf{[Calibration] 平均絶対誤差 (MAE)} & \textbf{1.0685} & 人間正解値との絶対スケール距離 \\
\textbf{[Calibration] 二乗平均平方根誤差 (RMSE)} & 1.2236 & 大きな乖離に対する感度 \\
\textbf{[分布] モデル予測期待値平均 $\pm$ SD} & 2.208 $\pm$ 0.870 & $E[V]$ の連続分布 \\
\textbf{[分布] 人間アノテーション平均 $\pm$ SD} & 2.978 $\pm$ 0.428 & 人間正解値の分布 \\
\textbf{[Collapse] 単一次元 Greedy=5 固定率} & \texttt{0.10\textbackslash{}\%} & 当該次元の最頻出力が中立に退化した割合 \\
\bottomrule
\end{tabular}
\end{center}


\textbf{表 5: ② Reader-Response Prediction (R) — Arousal}
\begin{center}
\small
\begin{tabular}{lcc}
\toprule
\textbf{評価軸 / 指標 (Metric)} & \textbf{値 (Value)} & \textbf{学術的意味・備考} \\
\midrule
\textbf{[Alignment] 連続相関 ($r$)} & \textbf{0.0345} & $p = 0.275 $ (刺激間パターンの追従度) \\
\textbf{[Alignment] 順位相関 ($\rho$)} & 0.0148 & $p = 0.639 $ (単調増加性) \\
\textbf{[Alignment] Greedy相関 ($r_{greedy}$)} & 0.0400 & 最尤トークン出力ベースの一致度 \\
\textbf{[Calibration] 平均絶対誤差 (MAE)} & \textbf{0.7662} & 人間正解値との絶対スケール距離 \\
\textbf{[Calibration] 二乗平均平方根誤差 (RMSE)} & 0.9410 & 大きな乖離に対する感度 \\
\textbf{[分布] モデル予測期待値平均 $\pm$ SD} & 2.546 $\pm$ 0.741 & $E[A]$ の連続分布 \\
\textbf{[分布] 人間アノテーション平均 $\pm$ SD} & 3.050 $\pm$ 0.313 & 人間正解値の分布 \\
\textbf{[Collapse] 単一次元 Greedy=5 固定率} & \texttt{0.40\textbackslash{}\%} & 当該次元の最頻出力が中立に退化した割合 \\
\bottomrule
\end{tabular}
\end{center}


\textbf{表 6: ② Reader-Response Prediction (R) — Dominance}
\begin{center}
\small
\begin{tabular}{lcc}
\toprule
\textbf{評価軸 / 指標 (Metric)} & \textbf{値 (Value)} & \textbf{学術的意味・備考} \\
\midrule
\textbf{[Alignment] 連続相関 ($r$)} & \textbf{-0.0079} & $p = 0.802 $ (刺激間パターンの追従度) \\
\textbf{[Alignment] 順位相関 ($\rho$)} & -0.0249 & $p = 0.431 $ (単調増加性) \\
\textbf{[Alignment] Greedy相関 ($r_{greedy}$)} & 0.0037 & 最尤トークン出力ベースの一致度 \\
\textbf{[Calibration] 平均絶対誤差 (MAE)} & \textbf{0.6808} & 人間正解値との絶対スケール距離 \\
\textbf{[Calibration] 二乗平均平方根誤差 (RMSE)} & 0.8426 & 大きな乖離に対する感度 \\
\textbf{[分布] モデル予測期待値平均 $\pm$ SD} & 2.503 $\pm$ 0.578 & $E[D]$ の連続分布 \\
\textbf{[分布] 人間アノテーション平均 $\pm$ SD} & 3.043 $\pm$ 0.287 & 人間正解値の分布 \\
\textbf{[Collapse] 単一次元 Greedy=5 固定率} & \texttt{0.30\textbackslash{}\%} & 当該次元の最頻出力が中立に退化した割合 \\
\bottomrule
\end{tabular}
\end{center}


\subparagraph{③ Self-Report (S)}
\textit{評価基準: human\_reader VAD}

\textbf{表 7: ③ Self-Report (S) — Valence}
\begin{center}
\small
\begin{tabular}{lcc}
\toprule
\textbf{評価軸 / 指標 (Metric)} & \textbf{値 (Value)} & \textbf{学術的意味・備考} \\
\midrule
\textbf{[Alignment] 連続相関 ($r$)} & \textbf{0.0548} & $p = 0.083 $ (刺激間パターンの追従度) \\
\textbf{[Alignment] 順位相関 ($\rho$)} & 0.0880 & $p = 0.005 $ (単調増加性) \\
\textbf{[Alignment] Greedy相関 ($r_{greedy}$)} & 0.0117 & 最尤トークン出力ベースの一致度 \\
\textbf{[Calibration] 平均絶対誤差 (MAE)} & \textbf{1.3886} & 人間正解値との絶対スケール距離 \\
\textbf{[Calibration] 二乗平均平方根誤差 (RMSE)} & 1.5071 & 大きな乖離に対する感度 \\
\textbf{[分布] モデル予測期待値平均 $\pm$ SD} & 1.754 $\pm$ 0.792 & $E[V]$ の連続分布 \\
\textbf{[分布] 人間アノテーション平均 $\pm$ SD} & 2.978 $\pm$ 0.428 & 人間正解値の分布 \\
\textbf{[Collapse] 単一次元 Greedy=5 固定率} & \texttt{0.00\textbackslash{}\%} & 当該次元の最頻出力が中立に退化した割合 \\
\bottomrule
\end{tabular}
\end{center}


\textbf{表 8: ③ Self-Report (S) — Arousal}
\begin{center}
\small
\begin{tabular}{lcc}
\toprule
\textbf{評価軸 / 指標 (Metric)} & \textbf{値 (Value)} & \textbf{学術的意味・備考} \\
\midrule
\textbf{[Alignment] 連続相関 ($r$)} & \textbf{0.0044} & $p = 0.889 $ (刺激間パターンの追従度) \\
\textbf{[Alignment] 順位相関 ($\rho$)} & -0.0089 & $p = 0.778 $ (単調増加性) \\
\textbf{[Alignment] Greedy相関 ($r_{greedy}$)} & 0.0051 & 最尤トークン出力ベースの一致度 \\
\textbf{[Calibration] 平均絶対誤差 (MAE)} & \textbf{1.2103} & 人間正解値との絶対スケール距離 \\
\textbf{[Calibration] 二乗平均平方根誤差 (RMSE)} & 1.3571 & 大きな乖離に対する感度 \\
\textbf{[分布] モデル予測期待値平均 $\pm$ SD} & 1.935 $\pm$ 0.709 & $E[A]$ の連続分布 \\
\textbf{[分布] 人間アノテーション平均 $\pm$ SD} & 3.050 $\pm$ 0.313 & 人間正解値の分布 \\
\textbf{[Collapse] 単一次元 Greedy=5 固定率} & \texttt{0.30\textbackslash{}\%} & 当該次元の最頻出力が中立に退化した割合 \\
\bottomrule
\end{tabular}
\end{center}


\textbf{表 9: ③ Self-Report (S) — Dominance}
\begin{center}
\small
\begin{tabular}{lcc}
\toprule
\textbf{評価軸 / 指標 (Metric)} & \textbf{値 (Value)} & \textbf{学術的意味・備考} \\
\midrule
\textbf{[Alignment] 連続相関 ($r$)} & \textbf{0.0001} & $p = 0.999 $ (刺激間パターンの追従度) \\
\textbf{[Alignment] 順位相関 ($\rho$)} & 0.0061 & $p = 0.848 $ (単調増加性) \\
\textbf{[Alignment] Greedy相関 ($r_{greedy}$)} & 0.0118 & 最尤トークン出力ベースの一致度 \\
\textbf{[Calibration] 平均絶対誤差 (MAE)} & \textbf{1.1162} & 人間正解値との絶対スケール距離 \\
\textbf{[Calibration] 二乗平均平方根誤差 (RMSE)} & 1.2561 & 大きな乖離に対する感度 \\
\textbf{[分布] モデル予測期待値平均 $\pm$ SD} & 2.000 $\pm$ 0.639 & $E[D]$ の連続分布 \\
\textbf{[分布] 人間アノテーション平均 $\pm$ SD} & 3.043 $\pm$ 0.287 & 人間正解値の分布 \\
\textbf{[Collapse] 単一次元 Greedy=5 固定率} & \texttt{0.20\textbackslash{}\%} & 当該次元の最頻出力が中立に退化した割合 \\
\bottomrule
\end{tabular}
\end{center}


\vspace{0.5em}\hrule\vspace{0.5em}

\paragraph{■ モデル: \texttt{gemma2\textbackslash{}\_2b\textbackslash{}\_instruct} (Instruct, N=1000)}

\subparagraph{【サマリー】 3×3 相関 \& 誤差マトリクス}
\begin{center}
\small
\begin{tabular}{lcccc}
\toprule
\textbf{タスク \ 感情次元} & \textbf{Valence ($V$)} & \textbf{Arousal ($A$)} & \textbf{Dominance ($D$)} & \textbf{完全中立率 $(5,5,5)$} \\
\midrule
\textbf{① Writer-State Estimation (W)} & \textbf{$r=0.319$} (MAE 0.64) & \textbf{$r=0.239$} (MAE 0.51) & \textbf{$r=0.098$} (MAE 0.50) & \texttt{0.0\textbackslash{}\%} \\
\textbf{② Reader-Response Prediction (R)} & \textbf{$r=0.247$} (MAE 0.60) & \textbf{$r=0.310$} (MAE 0.47) & \textbf{$r=0.025$} (MAE 0.59) & \texttt{0.0\textbackslash{}\%} \\
\textbf{③ Self-Report (S)} & \textbf{$r=0.379$} (MAE 0.58) & \textbf{$r=0.342$} (MAE 0.48) & \textbf{$r=0.069$} (MAE 0.49) & \texttt{0.0\textbackslash{}\%} \\
\bottomrule
\end{tabular}
\end{center}


\begin{itemize}
\item \textbf{内部結合度 (Internal Coupling $\text{Corr}(R, S)$)}: Valence $r=0.893$, Arousal $r=0.900$, Dominance $r=0.750
\end{itemize}

\subparagraph{【9個の詳細表】 タスク別・次元別詳細メトリクス}

\subparagraph{① Writer-State Estimation (W)}
\textit{評価基準: human\_writer VAD}

\textbf{表 1: ① Writer-State Estimation (W) — Valence}
\begin{center}
\small
\begin{tabular}{lcc}
\toprule
\textbf{評価軸 / 指標 (Metric)} & \textbf{値 (Value)} & \textbf{学術的意味・備考} \\
\midrule
\textbf{[Alignment] 連続相関 ($r$)} & \textbf{0.3190} & $p = <.001 $ (刺激間パターンの追従度) \\
\textbf{[Alignment] 順位相関 ($\rho$)} & 0.3354 & $p = <.001 $ (単調増加性) \\
\textbf{[Alignment] Greedy相関 ($r_{greedy}$)} & 0.2923 & 最尤トークン出力ベースの一致度 \\
\textbf{[Calibration] 平均絶対誤差 (MAE)} & \textbf{0.6370} & 人間正解値との絶対スケール距離 \\
\textbf{[Calibration] 二乗平均平方根誤差 (RMSE)} & 0.7734 & 大きな乖離に対する感度 \\
\textbf{[分布] モデル予測期待値平均 $\pm$ SD} & 3.211 $\pm$ 0.772 & $E[V]$ の連続分布 \\
\textbf{[分布] 人間アノテーション平均 $\pm$ SD} & 2.976 $\pm$ 0.335 & 人間正解値の分布 \\
\textbf{[Collapse] 単一次元 Greedy=5 固定率} & \texttt{18.50\textbackslash{}\%} & 当該次元の最頻出力が中立に退化した割合 \\
\bottomrule
\end{tabular}
\end{center}


\textbf{表 2: ① Writer-State Estimation (W) — Arousal}
\begin{center}
\small
\begin{tabular}{lcc}
\toprule
\textbf{評価軸 / 指標 (Metric)} & \textbf{値 (Value)} & \textbf{学術的意味・備考} \\
\midrule
\textbf{[Alignment] 連続相関 ($r$)} & \textbf{0.2387} & $p = <.001 $ (刺激間パターンの追従度) \\
\textbf{[Alignment] 順位相関 ($\rho$)} & 0.1958 & $p = <.001 $ (単調増加性) \\
\textbf{[Alignment] Greedy相関 ($r_{greedy}$)} & 0.2005 & 最尤トークン出力ベースの一致度 \\
\textbf{[Calibration] 平均絶対誤差 (MAE)} & \textbf{0.5071} & 人間正解値との絶対スケール距離 \\
\textbf{[Calibration] 二乗平均平方根誤差 (RMSE)} & 0.6236 & 大きな乖離に対する感度 \\
\textbf{[分布] モデル予測期待値平均 $\pm$ SD} & 3.402 $\pm$ 0.465 & $E[A]$ の連続分布 \\
\textbf{[分布] 人間アノテーション平均 $\pm$ SD} & 3.017 $\pm$ 0.305 & 人間正解値の分布 \\
\textbf{[Collapse] 単一次元 Greedy=5 固定率} & \texttt{24.80\textbackslash{}\%} & 当該次元の最頻出力が中立に退化した割合 \\
\bottomrule
\end{tabular}
\end{center}


\textbf{表 3: ① Writer-State Estimation (W) — Dominance}
\begin{center}
\small
\begin{tabular}{lcc}
\toprule
\textbf{評価軸 / 指標 (Metric)} & \textbf{値 (Value)} & \textbf{学術的意味・備考} \\
\midrule
\textbf{[Alignment] 連続相関 ($r$)} & \textbf{0.0980} & $p = 0.002 $ (刺激間パターンの追従度) \\
\textbf{[Alignment] 順位相関 ($\rho$)} & 0.0364 & $p = 0.251 $ (単調増加性) \\
\textbf{[Alignment] Greedy相関 ($r_{greedy}$)} & 0.0390 & 最尤トークン出力ベースの一致度 \\
\textbf{[Calibration] 平均絶対誤差 (MAE)} & \textbf{0.5014} & 人間正解値との絶対スケール距離 \\
\textbf{[Calibration] 二乗平均平方根誤差 (RMSE)} & 0.5997 & 大きな乖離に対する感度 \\
\textbf{[分布] モデル予測期待値平均 $\pm$ SD} & 2.657 $\pm$ 0.354 & $E[D]$ の連続分布 \\
\textbf{[分布] 人間アノテーション平均 $\pm$ SD} & 3.092 $\pm$ 0.251 & 人間正解値の分布 \\
\textbf{[Collapse] 単一次元 Greedy=5 固定率} & \texttt{45.20\textbackslash{}\%} & 当該次元の最頻出力が中立に退化した割合 \\
\bottomrule
\end{tabular}
\end{center}


\subparagraph{② Reader-Response Prediction (R)}
\textit{評価基準: human\_reader VAD}

\textbf{表 4: ② Reader-Response Prediction (R) — Valence}
\begin{center}
\small
\begin{tabular}{lcc}
\toprule
\textbf{評価軸 / 指標 (Metric)} & \textbf{値 (Value)} & \textbf{学術的意味・備考} \\
\midrule
\textbf{[Alignment] 連続相関 ($r$)} & \textbf{0.2473} & $p = <.001 $ (刺激間パターンの追従度) \\
\textbf{[Alignment] 順位相関 ($\rho$)} & 0.2995 & $p = <.001 $ (単調増加性) \\
\textbf{[Alignment] Greedy相関 ($r_{greedy}$)} & 0.2118 & 最尤トークン出力ベースの一致度 \\
\textbf{[Calibration] 平均絶対誤差 (MAE)} & \textbf{0.5994} & 人間正解値との絶対スケール距離 \\
\textbf{[Calibration] 二乗平均平方根誤差 (RMSE)} & 0.7678 & 大きな乖離に対する感度 \\
\textbf{[分布] モデル予測期待値平均 $\pm$ SD} & 3.253 $\pm$ 0.691 & $E[V]$ の連続分布 \\
\textbf{[分布] 人間アノテーション平均 $\pm$ SD} & 2.978 $\pm$ 0.428 & 人間正解値の分布 \\
\textbf{[Collapse] 単一次元 Greedy=5 固定率} & \texttt{23.70\textbackslash{}\%} & 当該次元の最頻出力が中立に退化した割合 \\
\bottomrule
\end{tabular}
\end{center}


\textbf{表 5: ② Reader-Response Prediction (R) — Arousal}
\begin{center}
\small
\begin{tabular}{lcc}
\toprule
\textbf{評価軸 / 指標 (Metric)} & \textbf{値 (Value)} & \textbf{学術的意味・備考} \\
\midrule
\textbf{[Alignment] 連続相関 ($r$)} & \textbf{0.3104} & $p = <.001 $ (刺激間パターンの追従度) \\
\textbf{[Alignment] 順位相関 ($\rho$)} & 0.3137 & $p = <.001 $ (単調増加性) \\
\textbf{[Alignment] Greedy相関 ($r_{greedy}$)} & 0.2466 & 最尤トークン出力ベースの一致度 \\
\textbf{[Calibration] 平均絶対誤差 (MAE)} & \textbf{0.4686} & 人間正解値との絶対スケール距離 \\
\textbf{[Calibration] 二乗平均平方根誤差 (RMSE)} & 0.5797 & 大きな乖離に対する感度 \\
\textbf{[分布] モデル予測期待値平均 $\pm$ SD} & 3.239 $\pm$ 0.558 & $E[A]$ の連続分布 \\
\textbf{[分布] 人間アノテーション平均 $\pm$ SD} & 3.050 $\pm$ 0.313 & 人間正解値の分布 \\
\textbf{[Collapse] 単一次元 Greedy=5 固定率} & \texttt{27.00\textbackslash{}\%} & 当該次元の最頻出力が中立に退化した割合 \\
\bottomrule
\end{tabular}
\end{center}


\textbf{表 6: ② Reader-Response Prediction (R) — Dominance}
\begin{center}
\small
\begin{tabular}{lcc}
\toprule
\textbf{評価軸 / 指標 (Metric)} & \textbf{値 (Value)} & \textbf{学術的意味・備考} \\
\midrule
\textbf{[Alignment] 連続相関 ($r$)} & \textbf{0.0250} & $p = 0.430 $ (刺激間パターンの追従度) \\
\textbf{[Alignment] 順位相関 ($\rho$)} & -0.0007 & $p = 0.981 $ (単調増加性) \\
\textbf{[Alignment] Greedy相関 ($r_{greedy}$)} & 0.0154 & 最尤トークン出力ベースの一致度 \\
\textbf{[Calibration] 平均絶対誤差 (MAE)} & \textbf{0.5890} & 人間正解値との絶対スケール距離 \\
\textbf{[Calibration] 二乗平均平方根誤差 (RMSE)} & 0.6991 & 大きな乖離に対する感度 \\
\textbf{[分布] モデル予測期待値平均 $\pm$ SD} & 2.576 $\pm$ 0.442 & $E[D]$ の連続分布 \\
\textbf{[分布] 人間アノテーション平均 $\pm$ SD} & 3.043 $\pm$ 0.287 & 人間正解値の分布 \\
\textbf{[Collapse] 単一次元 Greedy=5 固定率} & \texttt{33.20\textbackslash{}\%} & 当該次元の最頻出力が中立に退化した割合 \\
\bottomrule
\end{tabular}
\end{center}


\subparagraph{③ Self-Report (S)}
\textit{評価基準: human\_reader VAD}

\textbf{表 7: ③ Self-Report (S) — Valence}
\begin{center}
\small
\begin{tabular}{lcc}
\toprule
\textbf{評価軸 / 指標 (Metric)} & \textbf{値 (Value)} & \textbf{学術的意味・備考} \\
\midrule
\textbf{[Alignment] 連続相関 ($r$)} & \textbf{0.3794} & $p = <.001 $ (刺激間パターンの追従度) \\
\textbf{[Alignment] 順位相関 ($\rho$)} & 0.4321 & $p = <.001 $ (単調増加性) \\
\textbf{[Alignment] Greedy相関 ($r_{greedy}$)} & 0.3546 & 最尤トークン出力ベースの一致度 \\
\textbf{[Calibration] 平均絶対誤差 (MAE)} & \textbf{0.5770} & 人間正解値との絶対スケール距離 \\
\textbf{[Calibration] 二乗平均平方根誤差 (RMSE)} & 0.7347 & 大きな乖離に対する感度 \\
\textbf{[分布] モデル予測期待値平均 $\pm$ SD} & 3.329 $\pm$ 0.672 & $E[V]$ の連続分布 \\
\textbf{[分布] 人間アノテーション平均 $\pm$ SD} & 2.978 $\pm$ 0.428 & 人間正解値の分布 \\
\textbf{[Collapse] 単一次元 Greedy=5 固定率} & \texttt{21.90\textbackslash{}\%} & 当該次元の最頻出力が中立に退化した割合 \\
\bottomrule
\end{tabular}
\end{center}


\textbf{表 8: ③ Self-Report (S) — Arousal}
\begin{center}
\small
\begin{tabular}{lcc}
\toprule
\textbf{評価軸 / 指標 (Metric)} & \textbf{値 (Value)} & \textbf{学術的意味・備考} \\
\midrule
\textbf{[Alignment] 連続相関 ($r$)} & \textbf{0.3416} & $p = <.001 $ (刺激間パターンの追従度) \\
\textbf{[Alignment] 順位相関 ($\rho$)} & 0.3255 & $p = <.001 $ (単調増加性) \\
\textbf{[Alignment] Greedy相関 ($r_{greedy}$)} & 0.2941 & 最尤トークン出力ベースの一致度 \\
\textbf{[Calibration] 平均絶対誤差 (MAE)} & \textbf{0.4766} & 人間正解値との絶対スケール距離 \\
\textbf{[Calibration] 二乗平均平方根誤差 (RMSE)} & 0.5819 & 大きな乖離に対する感度 \\
\textbf{[分布] モデル予測期待値平均 $\pm$ SD} & 3.434 $\pm$ 0.431 & $E[A]$ の連続分布 \\
\textbf{[分布] 人間アノテーション平均 $\pm$ SD} & 3.050 $\pm$ 0.313 & 人間正解値の分布 \\
\textbf{[Collapse] 単一次元 Greedy=5 固定率} & \texttt{25.10\textbackslash{}\%} & 当該次元の最頻出力が中立に退化した割合 \\
\bottomrule
\end{tabular}
\end{center}


\textbf{表 9: ③ Self-Report (S) — Dominance}
\begin{center}
\small
\begin{tabular}{lcc}
\toprule
\textbf{評価軸 / 指標 (Metric)} & \textbf{値 (Value)} & \textbf{学術的意味・備考} \\
\midrule
\textbf{[Alignment] 連続相関 ($r$)} & \textbf{0.0693} & $p = 0.028 $ (刺激間パターンの追従度) \\
\textbf{[Alignment] 順位相関 ($\rho$)} & 0.0490 & $p = 0.121 $ (単調増加性) \\
\textbf{[Alignment] Greedy相関 ($r_{greedy}$)} & 0.0126 & 最尤トークン出力ベースの一致度 \\
\textbf{[Calibration] 平均絶対誤差 (MAE)} & \textbf{0.4941} & 人間正解値との絶対スケール距離 \\
\textbf{[Calibration] 二乗平均平方根誤差 (RMSE)} & 0.6031 & 大きな乖離に対する感度 \\
\textbf{[分布] モデル予測期待値平均 $\pm$ SD} & 2.674 $\pm$ 0.401 & $E[D]$ の連続分布 \\
\textbf{[分布] 人間アノテーション平均 $\pm$ SD} & 3.043 $\pm$ 0.287 & 人間正解値の分布 \\
\textbf{[Collapse] 単一次元 Greedy=5 固定率} & \texttt{37.00\textbackslash{}\%} & 当該次元の最頻出力が中立に退化した割合 \\
\bottomrule
\end{tabular}
\end{center}


\vspace{0.5em}\hrule\vspace{0.5em}

\paragraph{■ モデル: \texttt{llama3.2\textbackslash{}\_1b\textbackslash{}\_base} (Base, N=1000)}

\subparagraph{【サマリー】 3×3 相関 \& 誤差マトリクス}
\begin{center}
\small
\begin{tabular}{lcccc}
\toprule
\textbf{タスク \ 感情次元} & \textbf{Valence ($V$)} & \textbf{Arousal ($A$)} & \textbf{Dominance ($D$)} & \textbf{完全中立率 $(5,5,5)$} \\
\midrule
\textbf{① Writer-State Estimation (W)} & \textbf{$r=0.255$} (MAE 0.24) & \textbf{$r=-0.031$} (MAE 0.26) & \textbf{$r=-0.036$} (MAE 0.26) & \texttt{0.0\textbackslash{}\%} \\
\textbf{② Reader-Response Prediction (R)} & \textbf{$r=0.182$} (MAE 0.34) & \textbf{$r=0.108$} (MAE 0.29) & \textbf{$r=0.070$} (MAE 0.30) & \texttt{0.0\textbackslash{}\%} \\
\textbf{③ Self-Report (S)} & \textbf{$r=0.257$} (MAE 0.32) & \textbf{$r=0.048$} (MAE 0.27) & \textbf{$r=0.080$} (MAE 0.24) & \texttt{0.0\textbackslash{}\%} \\
\bottomrule
\end{tabular}
\end{center}


\begin{itemize}
\item \textbf{内部結合度 (Internal Coupling $\text{Corr}(R, S)$)}: Valence $r=0.907$, Arousal $r=0.923$, Dominance $r=0.886
\end{itemize}

\subparagraph{【9個の詳細表】 タスク別・次元別詳細メトリクス}

\subparagraph{① Writer-State Estimation (W)}
\textit{評価基準: human\_writer VAD}

\textbf{表 1: ① Writer-State Estimation (W) — Valence}
\begin{center}
\small
\begin{tabular}{lcc}
\toprule
\textbf{評価軸 / 指標 (Metric)} & \textbf{値 (Value)} & \textbf{学術的意味・備考} \\
\midrule
\textbf{[Alignment] 連続相関 ($r$)} & \textbf{0.2553} & $p = <.001 $ (刺激間パターンの追従度) \\
\textbf{[Alignment] 順位相関 ($\rho$)} & 0.3415 & $p = <.001 $ (単調増加性) \\
\textbf{[Alignment] Greedy相関 ($r_{greedy}$)} & 0.0175 & 最尤トークン出力ベースの一致度 \\
\textbf{[Calibration] 平均絶対誤差 (MAE)} & \textbf{0.2438} & 人間正解値との絶対スケール距離 \\
\textbf{[Calibration] 二乗平均平方根誤差 (RMSE)} & 0.3384 & 大きな乖離に対する感度 \\
\textbf{[分布] モデル予測期待値平均 $\pm$ SD} & 2.879 $\pm$ 0.100 & $E[V]$ の連続分布 \\
\textbf{[分布] 人間アノテーション平均 $\pm$ SD} & 2.976 $\pm$ 0.335 & 人間正解値の分布 \\
\textbf{[Collapse] 単一次元 Greedy=5 固定率} & \texttt{0.00\textbackslash{}\%} & 当該次元の最頻出力が中立に退化した割合 \\
\bottomrule
\end{tabular}
\end{center}


\textbf{表 2: ① Writer-State Estimation (W) — Arousal}
\begin{center}
\small
\begin{tabular}{lcc}
\toprule
\textbf{評価軸 / 指標 (Metric)} & \textbf{値 (Value)} & \textbf{学術的意味・備考} \\
\midrule
\textbf{[Alignment] 連続相関 ($r$)} & \textbf{-0.0307} & $p = 0.333 $ (刺激間パターンの追従度) \\
\textbf{[Alignment] 順位相関 ($\rho$)} & -0.0280 & $p = 0.377 $ (単調増加性) \\
\textbf{[Alignment] Greedy相関 ($r_{greedy}$)} & -0.0091 & 最尤トークン出力ベースの一致度 \\
\textbf{[Calibration] 平均絶対誤差 (MAE)} & \textbf{0.2566} & 人間正解値との絶対スケール距離 \\
\textbf{[Calibration] 二乗平均平方根誤差 (RMSE)} & 0.3381 & 大きな乖離に対する感度 \\
\textbf{[分布] モデル予測期待値平均 $\pm$ SD} & 2.905 $\pm$ 0.086 & $E[A]$ の連続分布 \\
\textbf{[分布] 人間アノテーション平均 $\pm$ SD} & 3.017 $\pm$ 0.305 & 人間正解値の分布 \\
\textbf{[Collapse] 単一次元 Greedy=5 固定率} & \texttt{0.00\textbackslash{}\%} & 当該次元の最頻出力が中立に退化した割合 \\
\bottomrule
\end{tabular}
\end{center}


\textbf{表 3: ① Writer-State Estimation (W) — Dominance}
\begin{center}
\small
\begin{tabular}{lcc}
\toprule
\textbf{評価軸 / 指標 (Metric)} & \textbf{値 (Value)} & \textbf{学術的意味・備考} \\
\midrule
\textbf{[Alignment] 連続相関 ($r$)} & \textbf{-0.0356} & $p = 0.261 $ (刺激間パターンの追従度) \\
\textbf{[Alignment] 順位相関 ($\rho$)} & -0.0247 & $p = 0.435 $ (単調増加性) \\
\textbf{[Alignment] Greedy相関 ($r_{greedy}$)} & 0.0108 & 最尤トークン出力ベースの一致度 \\
\textbf{[Calibration] 平均絶対誤差 (MAE)} & \textbf{0.2597} & 人間正解値との絶対スケール距離 \\
\textbf{[Calibration] 二乗平均平方根誤差 (RMSE)} & 0.3349 & 大きな乖離に対する感度 \\
\textbf{[分布] モデル予測期待値平均 $\pm$ SD} & 2.887 $\pm$ 0.076 & $E[D]$ の連続分布 \\
\textbf{[分布] 人間アノテーション平均 $\pm$ SD} & 3.092 $\pm$ 0.251 & 人間正解値の分布 \\
\textbf{[Collapse] 単一次元 Greedy=5 固定率} & \texttt{0.00\textbackslash{}\%} & 当該次元の最頻出力が中立に退化した割合 \\
\bottomrule
\end{tabular}
\end{center}


\subparagraph{② Reader-Response Prediction (R)}
\textit{評価基準: human\_reader VAD}

\textbf{表 4: ② Reader-Response Prediction (R) — Valence}
\begin{center}
\small
\begin{tabular}{lcc}
\toprule
\textbf{評価軸 / 指標 (Metric)} & \textbf{値 (Value)} & \textbf{学術的意味・備考} \\
\midrule
\textbf{[Alignment] 連続相関 ($r$)} & \textbf{0.1818} & $p = <.001 $ (刺激間パターンの追従度) \\
\textbf{[Alignment] 順位相関 ($\rho$)} & 0.2326 & $p = <.001 $ (単調増加性) \\
\textbf{[Alignment] Greedy相関 ($r_{greedy}$)} & N/A (std=0) & 最尤トークン出力ベースの一致度 \\
\textbf{[Calibration] 平均絶対誤差 (MAE)} & \textbf{0.3449} & 人間正解値との絶対スケール距離 \\
\textbf{[Calibration] 二乗平均平方根誤差 (RMSE)} & 0.4469 & 大きな乖離に対する感度 \\
\textbf{[分布] モデル予測期待値平均 $\pm$ SD} & 2.829 $\pm$ 0.090 & $E[V]$ の連続分布 \\
\textbf{[分布] 人間アノテーション平均 $\pm$ SD} & 2.978 $\pm$ 0.428 & 人間正解値の分布 \\
\textbf{[Collapse] 単一次元 Greedy=5 固定率} & \texttt{0.00\textbackslash{}\%} & 当該次元の最頻出力が中立に退化した割合 \\
\bottomrule
\end{tabular}
\end{center}


\textbf{表 5: ② Reader-Response Prediction (R) — Arousal}
\begin{center}
\small
\begin{tabular}{lcc}
\toprule
\textbf{評価軸 / 指標 (Metric)} & \textbf{値 (Value)} & \textbf{学術的意味・備考} \\
\midrule
\textbf{[Alignment] 連続相関 ($r$)} & \textbf{0.1080} & $p = <.001 $ (刺激間パターンの追従度) \\
\textbf{[Alignment] 順位相関 ($\rho$)} & 0.0716 & $p = 0.024 $ (単調増加性) \\
\textbf{[Alignment] Greedy相関 ($r_{greedy}$)} & N/A (std=0) & 最尤トークン出力ベースの一致度 \\
\textbf{[Calibration] 平均絶対誤差 (MAE)} & \textbf{0.2913} & 人間正解値との絶対スケール距離 \\
\textbf{[Calibration] 二乗平均平方根誤差 (RMSE)} & 0.3660 & 大きな乖離に対する感度 \\
\textbf{[分布] モデル予測期待値平均 $\pm$ SD} & 2.863 $\pm$ 0.079 & $E[A]$ の連続分布 \\
\textbf{[分布] 人間アノテーション平均 $\pm$ SD} & 3.050 $\pm$ 0.313 & 人間正解値の分布 \\
\textbf{[Collapse] 単一次元 Greedy=5 固定率} & \texttt{0.00\textbackslash{}\%} & 当該次元の最頻出力が中立に退化した割合 \\
\bottomrule
\end{tabular}
\end{center}


\textbf{表 6: ② Reader-Response Prediction (R) — Dominance}
\begin{center}
\small
\begin{tabular}{lcc}
\toprule
\textbf{評価軸 / 指標 (Metric)} & \textbf{値 (Value)} & \textbf{学術的意味・備考} \\
\midrule
\textbf{[Alignment] 連続相関 ($r$)} & \textbf{0.0696} & $p = 0.028 $ (刺激間パターンの追従度) \\
\textbf{[Alignment] 順位相関 ($\rho$)} & 0.0587 & $p = 0.064 $ (単調増加性) \\
\textbf{[Alignment] Greedy相関 ($r_{greedy}$)} & N/A (std=0) & 最尤トークン出力ベースの一致度 \\
\textbf{[Calibration] 平均絶対誤差 (MAE)} & \textbf{0.2983} & 人間正解値との絶対スケール距離 \\
\textbf{[Calibration] 二乗平均平方根誤差 (RMSE)} & 0.3718 & 大きな乖離に対する感度 \\
\textbf{[分布] モデル予測期待値平均 $\pm$ SD} & 2.811 $\pm$ 0.071 & $E[D]$ の連続分布 \\
\textbf{[分布] 人間アノテーション平均 $\pm$ SD} & 3.043 $\pm$ 0.287 & 人間正解値の分布 \\
\textbf{[Collapse] 単一次元 Greedy=5 固定率} & \texttt{0.00\textbackslash{}\%} & 当該次元の最頻出力が中立に退化した割合 \\
\bottomrule
\end{tabular}
\end{center}


\subparagraph{③ Self-Report (S)}
\textit{評価基準: human\_reader VAD}

\textbf{表 7: ③ Self-Report (S) — Valence}
\begin{center}
\small
\begin{tabular}{lcc}
\toprule
\textbf{評価軸 / 指標 (Metric)} & \textbf{値 (Value)} & \textbf{学術的意味・備考} \\
\midrule
\textbf{[Alignment] 連続相関 ($r$)} & \textbf{0.2567} & $p = <.001 $ (刺激間パターンの追従度) \\
\textbf{[Alignment] 順位相関 ($\rho$)} & 0.3168 & $p = <.001 $ (単調増加性) \\
\textbf{[Alignment] Greedy相関 ($r_{greedy}$)} & 0.0459 & 最尤トークン出力ベースの一致度 \\
\textbf{[Calibration] 平均絶対誤差 (MAE)} & \textbf{0.3209} & 人間正解値との絶対スケール距離 \\
\textbf{[Calibration] 二乗平均平方根誤差 (RMSE)} & 0.4295 & 大きな乖離に対する感度 \\
\textbf{[分布] モデル予測期待値平均 $\pm$ SD} & 2.864 $\pm$ 0.106 & $E[V]$ の連続分布 \\
\textbf{[分布] 人間アノテーション平均 $\pm$ SD} & 2.978 $\pm$ 0.428 & 人間正解値の分布 \\
\textbf{[Collapse] 単一次元 Greedy=5 固定率} & \texttt{0.00\textbackslash{}\%} & 当該次元の最頻出力が中立に退化した割合 \\
\bottomrule
\end{tabular}
\end{center}


\textbf{表 8: ③ Self-Report (S) — Arousal}
\begin{center}
\small
\begin{tabular}{lcc}
\toprule
\textbf{評価軸 / 指標 (Metric)} & \textbf{値 (Value)} & \textbf{学術的意味・備考} \\
\midrule
\textbf{[Alignment] 連続相関 ($r$)} & \textbf{0.0479} & $p = 0.130 $ (刺激間パターンの追従度) \\
\textbf{[Alignment] 順位相関 ($\rho$)} & 0.0263 & $p = 0.406 $ (単調増加性) \\
\textbf{[Alignment] Greedy相関 ($r_{greedy}$)} & 0.0151 & 最尤トークン出力ベースの一致度 \\
\textbf{[Calibration] 平均絶対誤差 (MAE)} & \textbf{0.2729} & 人間正解値との絶対スケール距離 \\
\textbf{[Calibration] 二乗平均平方根誤差 (RMSE)} & 0.3514 & 大きな乖離に対する感度 \\
\textbf{[分布] モデル予測期待値平均 $\pm$ SD} & 2.908 $\pm$ 0.088 & $E[A]$ の連続分布 \\
\textbf{[分布] 人間アノテーション平均 $\pm$ SD} & 3.050 $\pm$ 0.313 & 人間正解値の分布 \\
\textbf{[Collapse] 単一次元 Greedy=5 固定率} & \texttt{0.00\textbackslash{}\%} & 当該次元の最頻出力が中立に退化した割合 \\
\bottomrule
\end{tabular}
\end{center}


\textbf{表 9: ③ Self-Report (S) — Dominance}
\begin{center}
\small
\begin{tabular}{lcc}
\toprule
\textbf{評価軸 / 指標 (Metric)} & \textbf{値 (Value)} & \textbf{学術的意味・備考} \\
\midrule
\textbf{[Alignment] 連続相関 ($r$)} & \textbf{0.0798} & $p = 0.012 $ (刺激間パターンの追従度) \\
\textbf{[Alignment] 順位相関 ($\rho$)} & 0.0729 & $p = 0.021 $ (単調増加性) \\
\textbf{[Alignment] Greedy相関 ($r_{greedy}$)} & 0.0173 & 最尤トークン出力ベースの一致度 \\
\textbf{[Calibration] 平均絶対誤差 (MAE)} & \textbf{0.2414} & 人間正解値との絶対スケール距離 \\
\textbf{[Calibration] 二乗平均平方根誤差 (RMSE)} & 0.3212 & 大きな乖離に対する感度 \\
\textbf{[分布] モデル予測期待値平均 $\pm$ SD} & 2.906 $\pm$ 0.074 & $E[D]$ の連続分布 \\
\textbf{[分布] 人間アノテーション平均 $\pm$ SD} & 3.043 $\pm$ 0.287 & 人間正解値の分布 \\
\textbf{[Collapse] 単一次元 Greedy=5 固定率} & \texttt{0.00\textbackslash{}\%} & 当該次元の最頻出力が中立に退化した割合 \\
\bottomrule
\end{tabular}
\end{center}


\vspace{0.5em}\hrule\vspace{0.5em}

\paragraph{■ モデル: \texttt{llama3.2\textbackslash{}\_1b\textbackslash{}\_instruct} (Instruct, N=1000)}

\subparagraph{【サマリー】 3×3 相関 \& 誤差マトリクス}
\begin{center}
\small
\begin{tabular}{lcccc}
\toprule
\textbf{タスク \ 感情次元} & \textbf{Valence ($V$)} & \textbf{Arousal ($A$)} & \textbf{Dominance ($D$)} & \textbf{完全中立率 $(5,5,5)$} \\
\midrule
\textbf{① Writer-State Estimation (W)} & \textbf{$r=0.519$} (MAE 0.32) & \textbf{$r=0.157$} (MAE 0.49) & \textbf{$r=0.081$} (MAE 0.38) & \texttt{0.0\textbackslash{}\%} \\
\textbf{② Reader-Response Prediction (R)} & \textbf{$r=0.443$} (MAE 0.45) & \textbf{$r=0.263$} (MAE 0.58) & \textbf{$r=0.153$} (MAE 0.41) & \texttt{0.0\textbackslash{}\%} \\
\textbf{③ Self-Report (S)} & \textbf{$r=0.463$} (MAE 0.41) & \textbf{$r=0.254$} (MAE 0.55) & \textbf{$r=0.162$} (MAE 0.41) & \texttt{0.0\textbackslash{}\%} \\
\bottomrule
\end{tabular}
\end{center}


\begin{itemize}
\item \textbf{内部結合度 (Internal Coupling $\text{Corr}(R, S)$)}: Valence $r=0.857$, Arousal $r=0.868$, Dominance $r=0.934
\end{itemize}

\subparagraph{【9個の詳細表】 タスク別・次元別詳細メトリクス}

\subparagraph{① Writer-State Estimation (W)}
\textit{評価基準: human\_writer VAD}

\textbf{表 1: ① Writer-State Estimation (W) — Valence}
\begin{center}
\small
\begin{tabular}{lcc}
\toprule
\textbf{評価軸 / 指標 (Metric)} & \textbf{値 (Value)} & \textbf{学術的意味・備考} \\
\midrule
\textbf{[Alignment] 連続相関 ($r$)} & \textbf{0.5192} & $p = <.001 $ (刺激間パターンの追従度) \\
\textbf{[Alignment] 順位相関 ($\rho$)} & 0.5225 & $p = <.001 $ (単調増加性) \\
\textbf{[Alignment] Greedy相関 ($r_{greedy}$)} & 0.3961 & 最尤トークン出力ベースの一致度 \\
\textbf{[Calibration] 平均絶対誤差 (MAE)} & \textbf{0.3190} & 人間正解値との絶対スケール距離 \\
\textbf{[Calibration] 二乗平均平方根誤差 (RMSE)} & 0.3988 & 大きな乖離に対する感度 \\
\textbf{[分布] モデル予測期待値平均 $\pm$ SD} & 2.764 $\pm$ 0.353 & $E[V]$ の連続分布 \\
\textbf{[分布] 人間アノテーション平均 $\pm$ SD} & 2.976 $\pm$ 0.335 & 人間正解値の分布 \\
\textbf{[Collapse] 単一次元 Greedy=5 固定率} & \texttt{0.10\textbackslash{}\%} & 当該次元の最頻出力が中立に退化した割合 \\
\bottomrule
\end{tabular}
\end{center}


\textbf{表 2: ① Writer-State Estimation (W) — Arousal}
\begin{center}
\small
\begin{tabular}{lcc}
\toprule
\textbf{評価軸 / 指標 (Metric)} & \textbf{値 (Value)} & \textbf{学術的意味・備考} \\
\midrule
\textbf{[Alignment] 連続相関 ($r$)} & \textbf{0.1574} & $p = <.001 $ (刺激間パターンの追従度) \\
\textbf{[Alignment] 順位相関 ($\rho$)} & 0.1266 & $p = <.001 $ (単調増加性) \\
\textbf{[Alignment] Greedy相関 ($r_{greedy}$)} & 0.1293 & 最尤トークン出力ベースの一致度 \\
\textbf{[Calibration] 平均絶対誤差 (MAE)} & \textbf{0.4947} & 人間正解値との絶対スケール距離 \\
\textbf{[Calibration] 二乗平均平方根誤差 (RMSE)} & 0.5821 & 大きな乖離に対する感度 \\
\textbf{[分布] モデル予測期待値平均 $\pm$ SD} & 2.572 $\pm$ 0.272 & $E[A]$ の連続分布 \\
\textbf{[分布] 人間アノテーション平均 $\pm$ SD} & 3.017 $\pm$ 0.305 & 人間正解値の分布 \\
\textbf{[Collapse] 単一次元 Greedy=5 固定率} & \texttt{0.10\textbackslash{}\%} & 当該次元の最頻出力が中立に退化した割合 \\
\bottomrule
\end{tabular}
\end{center}


\textbf{表 3: ① Writer-State Estimation (W) — Dominance}
\begin{center}
\small
\begin{tabular}{lcc}
\toprule
\textbf{評価軸 / 指標 (Metric)} & \textbf{値 (Value)} & \textbf{学術的意味・備考} \\
\midrule
\textbf{[Alignment] 連続相関 ($r$)} & \textbf{0.0812} & $p = 0.010 $ (刺激間パターンの追従度) \\
\textbf{[Alignment] 順位相関 ($\rho$)} & 0.0790 & $p = 0.012 $ (単調増加性) \\
\textbf{[Alignment] Greedy相関 ($r_{greedy}$)} & 0.0139 & 最尤トークン出力ベースの一致度 \\
\textbf{[Calibration] 平均絶対誤差 (MAE)} & \textbf{0.3808} & 人間正解値との絶対スケール距離 \\
\textbf{[Calibration] 二乗平均平方根誤差 (RMSE)} & 0.4630 & 大きな乖離に対する感度 \\
\textbf{[分布] モデル予測期待値平均 $\pm$ SD} & 2.784 $\pm$ 0.259 & $E[D]$ の連続分布 \\
\textbf{[分布] 人間アノテーション平均 $\pm$ SD} & 3.092 $\pm$ 0.251 & 人間正解値の分布 \\
\textbf{[Collapse] 単一次元 Greedy=5 固定率} & \texttt{0.40\textbackslash{}\%} & 当該次元の最頻出力が中立に退化した割合 \\
\bottomrule
\end{tabular}
\end{center}


\subparagraph{② Reader-Response Prediction (R)}
\textit{評価基準: human\_reader VAD}

\textbf{表 4: ② Reader-Response Prediction (R) — Valence}
\begin{center}
\small
\begin{tabular}{lcc}
\toprule
\textbf{評価軸 / 指標 (Metric)} & \textbf{値 (Value)} & \textbf{学術的意味・備考} \\
\midrule
\textbf{[Alignment] 連続相関 ($r$)} & \textbf{0.4426} & $p = <.001 $ (刺激間パターンの追従度) \\
\textbf{[Alignment] 順位相関 ($\rho$)} & 0.4490 & $p = <.001 $ (単調増加性) \\
\textbf{[Alignment] Greedy相関 ($r_{greedy}$)} & 0.2922 & 最尤トークン出力ベースの一致度 \\
\textbf{[Calibration] 平均絶対誤差 (MAE)} & \textbf{0.4495} & 人間正解値との絶対スケール距離 \\
\textbf{[Calibration] 二乗平均平方根誤差 (RMSE)} & 0.5337 & 大きな乖離に対する感度 \\
\textbf{[分布] モデル予測期待値平均 $\pm$ SD} & 2.647 $\pm$ 0.356 & $E[V]$ の連続分布 \\
\textbf{[分布] 人間アノテーション平均 $\pm$ SD} & 2.978 $\pm$ 0.428 & 人間正解値の分布 \\
\textbf{[Collapse] 単一次元 Greedy=5 固定率} & \texttt{0.00\textbackslash{}\%} & 当該次元の最頻出力が中立に退化した割合 \\
\bottomrule
\end{tabular}
\end{center}


\textbf{表 5: ② Reader-Response Prediction (R) — Arousal}
\begin{center}
\small
\begin{tabular}{lcc}
\toprule
\textbf{評価軸 / 指標 (Metric)} & \textbf{値 (Value)} & \textbf{学術的意味・備考} \\
\midrule
\textbf{[Alignment] 連続相関 ($r$)} & \textbf{0.2627} & $p = <.001 $ (刺激間パターンの追従度) \\
\textbf{[Alignment] 順位相関 ($\rho$)} & 0.2238 & $p = <.001 $ (単調増加性) \\
\textbf{[Alignment] Greedy相関 ($r_{greedy}$)} & 0.1386 & 最尤トークン出力ベースの一致度 \\
\textbf{[Calibration] 平均絶対誤差 (MAE)} & \textbf{0.5847} & 人間正解値との絶対スケール距離 \\
\textbf{[Calibration] 二乗平均平方根誤差 (RMSE)} & 0.6635 & 大きな乖離に対する感度 \\
\textbf{[分布] モデル予測期待値平均 $\pm$ SD} & 2.487 $\pm$ 0.260 & $E[A]$ の連続分布 \\
\textbf{[分布] 人間アノテーション平均 $\pm$ SD} & 3.050 $\pm$ 0.313 & 人間正解値の分布 \\
\textbf{[Collapse] 単一次元 Greedy=5 固定率} & \texttt{3.30\textbackslash{}\%} & 当該次元の最頻出力が中立に退化した割合 \\
\bottomrule
\end{tabular}
\end{center}


\textbf{表 6: ② Reader-Response Prediction (R) — Dominance}
\begin{center}
\small
\begin{tabular}{lcc}
\toprule
\textbf{評価軸 / 指標 (Metric)} & \textbf{値 (Value)} & \textbf{学術的意味・備考} \\
\midrule
\textbf{[Alignment] 連続相関 ($r$)} & \textbf{0.1534} & $p = <.001 $ (刺激間パターンの追従度) \\
\textbf{[Alignment] 順位相関 ($\rho$)} & 0.1503 & $p = <.001 $ (単調増加性) \\
\textbf{[Alignment] Greedy相関 ($r_{greedy}$)} & 0.0109 & 最尤トークン出力ベースの一致度 \\
\textbf{[Calibration] 平均絶対誤差 (MAE)} & \textbf{0.4080} & 人間正解値との絶対スケール距離 \\
\textbf{[Calibration] 二乗平均平方根誤差 (RMSE)} & 0.4937 & 大きな乖離に対する感度 \\
\textbf{[分布] モデル予測期待値平均 $\pm$ SD} & 2.708 $\pm$ 0.271 & $E[D]$ の連続分布 \\
\textbf{[分布] 人間アノテーション平均 $\pm$ SD} & 3.043 $\pm$ 0.287 & 人間正解値の分布 \\
\textbf{[Collapse] 単一次元 Greedy=5 固定率} & \texttt{0.50\textbackslash{}\%} & 当該次元の最頻出力が中立に退化した割合 \\
\bottomrule
\end{tabular}
\end{center}


\subparagraph{③ Self-Report (S)}
\textit{評価基準: human\_reader VAD}

\textbf{表 7: ③ Self-Report (S) — Valence}
\begin{center}
\small
\begin{tabular}{lcc}
\toprule
\textbf{評価軸 / 指標 (Metric)} & \textbf{値 (Value)} & \textbf{学術的意味・備考} \\
\midrule
\textbf{[Alignment] 連続相関 ($r$)} & \textbf{0.4634} & $p = <.001 $ (刺激間パターンの追従度) \\
\textbf{[Alignment] 順位相関 ($\rho$)} & 0.4764 & $p = <.001 $ (単調増加性) \\
\textbf{[Alignment] Greedy相関 ($r_{greedy}$)} & 0.2823 & 最尤トークン出力ベースの一致度 \\
\textbf{[Calibration] 平均絶対誤差 (MAE)} & \textbf{0.4063} & 人間正解値との絶対スケール距離 \\
\textbf{[Calibration] 二乗平均平方根誤差 (RMSE)} & 0.4895 & 大きな乖離に対する感度 \\
\textbf{[分布] モデル予測期待値平均 $\pm$ SD} & 2.733 $\pm$ 0.387 & $E[V]$ の連続分布 \\
\textbf{[分布] 人間アノテーション平均 $\pm$ SD} & 2.978 $\pm$ 0.428 & 人間正解値の分布 \\
\textbf{[Collapse] 単一次元 Greedy=5 固定率} & \texttt{5.40\textbackslash{}\%} & 当該次元の最頻出力が中立に退化した割合 \\
\bottomrule
\end{tabular}
\end{center}


\textbf{表 8: ③ Self-Report (S) — Arousal}
\begin{center}
\small
\begin{tabular}{lcc}
\toprule
\textbf{評価軸 / 指標 (Metric)} & \textbf{値 (Value)} & \textbf{学術的意味・備考} \\
\midrule
\textbf{[Alignment] 連続相関 ($r$)} & \textbf{0.2541} & $p = <.001 $ (刺激間パターンの追従度) \\
\textbf{[Alignment] 順位相関 ($\rho$)} & 0.2198 & $p = <.001 $ (単調増加性) \\
\textbf{[Alignment] Greedy相関 ($r_{greedy}$)} & 0.1710 & 最尤トークン出力ベースの一致度 \\
\textbf{[Calibration] 平均絶対誤差 (MAE)} & \textbf{0.5468} & 人間正解値との絶対スケール距離 \\
\textbf{[Calibration] 二乗平均平方根誤差 (RMSE)} & 0.6239 & 大きな乖離に対する感度 \\
\textbf{[分布] モデル予測期待値平均 $\pm$ SD} & 2.531 $\pm$ 0.246 & $E[A]$ の連続分布 \\
\textbf{[分布] 人間アノテーション平均 $\pm$ SD} & 3.050 $\pm$ 0.313 & 人間正解値の分布 \\
\textbf{[Collapse] 単一次元 Greedy=5 固定率} & \texttt{0.60\textbackslash{}\%} & 当該次元の最頻出力が中立に退化した割合 \\
\bottomrule
\end{tabular}
\end{center}


\textbf{表 9: ③ Self-Report (S) — Dominance}
\begin{center}
\small
\begin{tabular}{lcc}
\toprule
\textbf{評価軸 / 指標 (Metric)} & \textbf{値 (Value)} & \textbf{学術的意味・備考} \\
\midrule
\textbf{[Alignment] 連続相関 ($r$)} & \textbf{0.1617} & $p = <.001 $ (刺激間パターンの追従度) \\
\textbf{[Alignment] 順位相関 ($\rho$)} & 0.1501 & $p = <.001 $ (単調増加性) \\
\textbf{[Alignment] Greedy相関 ($r_{greedy}$)} & 0.0627 & 最尤トークン出力ベースの一致度 \\
\textbf{[Calibration] 平均絶対誤差 (MAE)} & \textbf{0.4128} & 人間正解値との絶対スケール距離 \\
\textbf{[Calibration] 二乗平均平方根誤差 (RMSE)} & 0.5006 & 大きな乖離に対する感度 \\
\textbf{[分布] モデル予測期待値平均 $\pm$ SD} & 2.703 $\pm$ 0.280 & $E[D]$ の連続分布 \\
\textbf{[分布] 人間アノテーション平均 $\pm$ SD} & 3.043 $\pm$ 0.287 & 人間正解値の分布 \\
\textbf{[Collapse] 単一次元 Greedy=5 固定率} & \texttt{2.70\textbackslash{}\%} & 当該次元の最頻出力が中立に退化した割合 \\
\bottomrule
\end{tabular}
\end{center}


\vspace{0.5em}\hrule\vspace{0.5em}

\paragraph{■ モデル: \texttt{mistral7b\textbackslash{}\_v0.1\textbackslash{}\_base} (Base, N=1000)}

\subparagraph{【サマリー】 3×3 相関 \& 誤差マトリクス}
\begin{center}
\small
\begin{tabular}{lcccc}
\toprule
\textbf{タスク \ 感情次元} & \textbf{Valence ($V$)} & \textbf{Arousal ($A$)} & \textbf{Dominance ($D$)} & \textbf{完全中立率 $(5,5,5)$} \\
\midrule
\textbf{① Writer-State Estimation (W)} & \textbf{$r=0.605$} (MAE 0.28) & \textbf{$r=0.252$} (MAE 0.24) & \textbf{$r=0.202$} (MAE 0.30) & \texttt{36.2\textbackslash{}\%} \\
\textbf{② Reader-Response Prediction (R)} & \textbf{$r=0.641$} (MAE 0.29) & \textbf{$r=0.337$} (MAE 0.28) & \textbf{$r=0.171$} (MAE 0.34) & \texttt{42.8\textbackslash{}\%} \\
\textbf{③ Self-Report (S)} & \textbf{$r=0.661$} (MAE 0.25) & \textbf{$r=0.326$} (MAE 0.25) & \textbf{$r=0.248$} (MAE 0.28) & \texttt{32.4\textbackslash{}\%} \\
\bottomrule
\end{tabular}
\end{center}


\begin{itemize}
\item \textbf{内部結合度 (Internal Coupling $\text{Corr}(R, S)$)}: Valence $r=0.962$, Arousal $r=0.902$, Dominance $r=0.852
\end{itemize}

\subparagraph{【9個の詳細表】 タスク別・次元別詳細メトリクス}

\subparagraph{① Writer-State Estimation (W)}
\textit{評価基準: human\_writer VAD}

\textbf{表 1: ① Writer-State Estimation (W) — Valence}
\begin{center}
\small
\begin{tabular}{lcc}
\toprule
\textbf{評価軸 / 指標 (Metric)} & \textbf{値 (Value)} & \textbf{学術的意味・備考} \\
\midrule
\textbf{[Alignment] 連続相関 ($r$)} & \textbf{0.6050} & $p = <.001 $ (刺激間パターンの追従度) \\
\textbf{[Alignment] 順位相関 ($\rho$)} & 0.6319 & $p = <.001 $ (単調増加性) \\
\textbf{[Alignment] Greedy相関 ($r_{greedy}$)} & 0.3361 & 最尤トークン出力ベースの一致度 \\
\textbf{[Calibration] 平均絶対誤差 (MAE)} & \textbf{0.2834} & 人間正解値との絶対スケール距離 \\
\textbf{[Calibration] 二乗平均平方根誤差 (RMSE)} & 0.3552 & 大きな乖離に対する感度 \\
\textbf{[分布] モデル予測期待値平均 $\pm$ SD} & 2.803 $\pm$ 0.361 & $E[V]$ の連続分布 \\
\textbf{[分布] 人間アノテーション平均 $\pm$ SD} & 2.976 $\pm$ 0.335 & 人間正解値の分布 \\
\textbf{[Collapse] 単一次元 Greedy=5 固定率} & \texttt{36.30\textbackslash{}\%} & 当該次元の最頻出力が中立に退化した割合 \\
\bottomrule
\end{tabular}
\end{center}


\textbf{表 2: ① Writer-State Estimation (W) — Arousal}
\begin{center}
\small
\begin{tabular}{lcc}
\toprule
\textbf{評価軸 / 指標 (Metric)} & \textbf{値 (Value)} & \textbf{学術的意味・備考} \\
\midrule
\textbf{[Alignment] 連続相関 ($r$)} & \textbf{0.2524} & $p = <.001 $ (刺激間パターンの追従度) \\
\textbf{[Alignment] 順位相関 ($\rho$)} & 0.2030 & $p = <.001 $ (単調増加性) \\
\textbf{[Alignment] Greedy相関 ($r_{greedy}$)} & 0.0754 & 最尤トークン出力ベースの一致度 \\
\textbf{[Calibration] 平均絶対誤差 (MAE)} & \textbf{0.2407} & 人間正解値との絶対スケール距離 \\
\textbf{[Calibration] 二乗平均平方根誤差 (RMSE)} & 0.3081 & 大きな乖離に対する感度 \\
\textbf{[分布] モデル予測期待値平均 $\pm$ SD} & 2.983 $\pm$ 0.160 & $E[A]$ の連続分布 \\
\textbf{[分布] 人間アノテーション平均 $\pm$ SD} & 3.017 $\pm$ 0.305 & 人間正解値の分布 \\
\textbf{[Collapse] 単一次元 Greedy=5 固定率} & \texttt{36.30\textbackslash{}\%} & 当該次元の最頻出力が中立に退化した割合 \\
\bottomrule
\end{tabular}
\end{center}


\textbf{表 3: ① Writer-State Estimation (W) — Dominance}
\begin{center}
\small
\begin{tabular}{lcc}
\toprule
\textbf{評価軸 / 指標 (Metric)} & \textbf{値 (Value)} & \textbf{学術的意味・備考} \\
\midrule
\textbf{[Alignment] 連続相関 ($r$)} & \textbf{0.2022} & $p = <.001 $ (刺激間パターンの追従度) \\
\textbf{[Alignment] 順位相関 ($\rho$)} & 0.1768 & $p = <.001 $ (単調増加性) \\
\textbf{[Alignment] Greedy相関 ($r_{greedy}$)} & 0.1449 & 最尤トークン出力ベースの一致度 \\
\textbf{[Calibration] 平均絶対誤差 (MAE)} & \textbf{0.3000} & 人間正解値との絶対スケール距離 \\
\textbf{[Calibration] 二乗平均平方根誤差 (RMSE)} & 0.3661 & 大きな乖離に対する感度 \\
\textbf{[分布] モデル予測期待値平均 $\pm$ SD} & 2.827 $\pm$ 0.109 & $E[D]$ の連続分布 \\
\textbf{[分布] 人間アノテーション平均 $\pm$ SD} & 3.092 $\pm$ 0.251 & 人間正解値の分布 \\
\textbf{[Collapse] 単一次元 Greedy=5 固定率} & \texttt{36.20\textbackslash{}\%} & 当該次元の最頻出力が中立に退化した割合 \\
\bottomrule
\end{tabular}
\end{center}


\subparagraph{② Reader-Response Prediction (R)}
\textit{評価基準: human\_reader VAD}

\textbf{表 4: ② Reader-Response Prediction (R) — Valence}
\begin{center}
\small
\begin{tabular}{lcc}
\toprule
\textbf{評価軸 / 指標 (Metric)} & \textbf{値 (Value)} & \textbf{学術的意味・備考} \\
\midrule
\textbf{[Alignment] 連続相関 ($r$)} & \textbf{0.6410} & $p = <.001 $ (刺激間パターンの追従度) \\
\textbf{[Alignment] 順位相関 ($\rho$)} & 0.6610 & $p = <.001 $ (単調増加性) \\
\textbf{[Alignment] Greedy相関 ($r_{greedy}$)} & 0.3464 & 最尤トークン出力ベースの一致度 \\
\textbf{[Calibration] 平均絶対誤差 (MAE)} & \textbf{0.2898} & 人間正解値との絶対スケール距離 \\
\textbf{[Calibration] 二乗平均平方根誤差 (RMSE)} & 0.3723 & 大きな乖離に対する感度 \\
\textbf{[分布] モデル予測期待値平均 $\pm$ SD} & 2.859 $\pm$ 0.402 & $E[V]$ の連続分布 \\
\textbf{[分布] 人間アノテーション平均 $\pm$ SD} & 2.978 $\pm$ 0.428 & 人間正解値の分布 \\
\textbf{[Collapse] 単一次元 Greedy=5 固定率} & \texttt{42.90\textbackslash{}\%} & 当該次元の最頻出力が中立に退化した割合 \\
\bottomrule
\end{tabular}
\end{center}


\textbf{表 5: ② Reader-Response Prediction (R) — Arousal}
\begin{center}
\small
\begin{tabular}{lcc}
\toprule
\textbf{評価軸 / 指標 (Metric)} & \textbf{値 (Value)} & \textbf{学術的意味・備考} \\
\midrule
\textbf{[Alignment] 連続相関 ($r$)} & \textbf{0.3366} & $p = <.001 $ (刺激間パターンの追従度) \\
\textbf{[Alignment] 順位相関 ($\rho$)} & 0.3111 & $p = <.001 $ (単調増加性) \\
\textbf{[Alignment] Greedy相関 ($r_{greedy}$)} & -0.0102 & 最尤トークン出力ベースの一致度 \\
\textbf{[Calibration] 平均絶対誤差 (MAE)} & \textbf{0.2751} & 人間正解値との絶対スケール距離 \\
\textbf{[Calibration] 二乗平均平方根誤差 (RMSE)} & 0.3447 & 大きな乖離に対する感度 \\
\textbf{[分布] モデル予測期待値平均 $\pm$ SD} & 2.905 $\pm$ 0.209 & $E[A]$ の連続分布 \\
\textbf{[分布] 人間アノテーション平均 $\pm$ SD} & 3.050 $\pm$ 0.313 & 人間正解値の分布 \\
\textbf{[Collapse] 単一次元 Greedy=5 固定率} & \texttt{43.90\textbackslash{}\%} & 当該次元の最頻出力が中立に退化した割合 \\
\bottomrule
\end{tabular}
\end{center}


\textbf{表 6: ② Reader-Response Prediction (R) — Dominance}
\begin{center}
\small
\begin{tabular}{lcc}
\toprule
\textbf{評価軸 / 指標 (Metric)} & \textbf{値 (Value)} & \textbf{学術的意味・備考} \\
\midrule
\textbf{[Alignment] 連続相関 ($r$)} & \textbf{0.1712} & $p = <.001 $ (刺激間パターンの追従度) \\
\textbf{[Alignment] 順位相関 ($\rho$)} & 0.1427 & $p = <.001 $ (単調増加性) \\
\textbf{[Alignment] Greedy相関 ($r_{greedy}$)} & 0.1673 & 最尤トークン出力ベースの一致度 \\
\textbf{[Calibration] 平均絶対誤差 (MAE)} & \textbf{0.3404} & 人間正解値との絶対スケール距離 \\
\textbf{[Calibration] 二乗平均平方根誤差 (RMSE)} & 0.4074 & 大きな乖離に対する感度 \\
\textbf{[分布] モデル予測期待値平均 $\pm$ SD} & 2.756 $\pm$ 0.110 & $E[D]$ の連続分布 \\
\textbf{[分布] 人間アノテーション平均 $\pm$ SD} & 3.043 $\pm$ 0.287 & 人間正解値の分布 \\
\textbf{[Collapse] 単一次元 Greedy=5 固定率} & \texttt{43.60\textbackslash{}\%} & 当該次元の最頻出力が中立に退化した割合 \\
\bottomrule
\end{tabular}
\end{center}


\subparagraph{③ Self-Report (S)}
\textit{評価基準: human\_reader VAD}

\textbf{表 7: ③ Self-Report (S) — Valence}
\begin{center}
\small
\begin{tabular}{lcc}
\toprule
\textbf{評価軸 / 指標 (Metric)} & \textbf{値 (Value)} & \textbf{学術的意味・備考} \\
\midrule
\textbf{[Alignment] 連続相関 ($r$)} & \textbf{0.6609} & $p = <.001 $ (刺激間パターンの追従度) \\
\textbf{[Alignment] 順位相関 ($\rho$)} & 0.6652 & $p = <.001 $ (単調増加性) \\
\textbf{[Alignment] Greedy相関 ($r_{greedy}$)} & 0.3183 & 最尤トークン出力ベースの一致度 \\
\textbf{[Calibration] 平均絶対誤差 (MAE)} & \textbf{0.2501} & 人間正解値との絶対スケール距離 \\
\textbf{[Calibration] 二乗平均平方根誤差 (RMSE)} & 0.3300 & 大きな乖離に対する感度 \\
\textbf{[分布] モデル予測期待値平均 $\pm$ SD} & 2.922 $\pm$ 0.333 & $E[V]$ の連続分布 \\
\textbf{[分布] 人間アノテーション平均 $\pm$ SD} & 2.978 $\pm$ 0.428 & 人間正解値の分布 \\
\textbf{[Collapse] 単一次元 Greedy=5 固定率} & \texttt{32.40\textbackslash{}\%} & 当該次元の最頻出力が中立に退化した割合 \\
\bottomrule
\end{tabular}
\end{center}


\textbf{表 8: ③ Self-Report (S) — Arousal}
\begin{center}
\small
\begin{tabular}{lcc}
\toprule
\textbf{評価軸 / 指標 (Metric)} & \textbf{値 (Value)} & \textbf{学術的意味・備考} \\
\midrule
\textbf{[Alignment] 連続相関 ($r$)} & \textbf{0.3263} & $p = <.001 $ (刺激間パターンの追従度) \\
\textbf{[Alignment] 順位相関 ($\rho$)} & 0.2764 & $p = <.001 $ (単調増加性) \\
\textbf{[Alignment] Greedy相関 ($r_{greedy}$)} & -0.0331 & 最尤トークン出力ベースの一致度 \\
\textbf{[Calibration] 平均絶対誤差 (MAE)} & \textbf{0.2545} & 人間正解値との絶対スケール距離 \\
\textbf{[Calibration] 二乗平均平方根誤差 (RMSE)} & 0.3210 & 大きな乖離に対する感度 \\
\textbf{[分布] モデル予測期待値平均 $\pm$ SD} & 2.932 $\pm$ 0.141 & $E[A]$ の連続分布 \\
\textbf{[分布] 人間アノテーション平均 $\pm$ SD} & 3.050 $\pm$ 0.313 & 人間正解値の分布 \\
\textbf{[Collapse] 単一次元 Greedy=5 固定率} & \texttt{32.40\textbackslash{}\%} & 当該次元の最頻出力が中立に退化した割合 \\
\bottomrule
\end{tabular}
\end{center}


\textbf{表 9: ③ Self-Report (S) — Dominance}
\begin{center}
\small
\begin{tabular}{lcc}
\toprule
\textbf{評価軸 / 指標 (Metric)} & \textbf{値 (Value)} & \textbf{学術的意味・備考} \\
\midrule
\textbf{[Alignment] 連続相関 ($r$)} & \textbf{0.2477} & $p = <.001 $ (刺激間パターンの追従度) \\
\textbf{[Alignment] 順位相関 ($\rho$)} & 0.2244 & $p = <.001 $ (単調増加性) \\
\textbf{[Alignment] Greedy相関 ($r_{greedy}$)} & 0.1353 & 最尤トークン出力ベースの一致度 \\
\textbf{[Calibration] 平均絶対誤差 (MAE)} & \textbf{0.2842} & 人間正解値との絶対スケール距離 \\
\textbf{[Calibration] 二乗平均平方根誤差 (RMSE)} & 0.3550 & 大きな乖離に対する感度 \\
\textbf{[分布] モデル予測期待値平均 $\pm$ SD} & 2.824 $\pm$ 0.100 & $E[D]$ の連続分布 \\
\textbf{[分布] 人間アノテーション平均 $\pm$ SD} & 3.043 $\pm$ 0.287 & 人間正解値の分布 \\
\textbf{[Collapse] 単一次元 Greedy=5 固定率} & \texttt{32.40\textbackslash{}\%} & 当該次元の最頻出力が中立に退化した割合 \\
\bottomrule
\end{tabular}
\end{center}


\vspace{0.5em}\hrule\vspace{0.5em}

\paragraph{■ モデル: \texttt{mistral7b\textbackslash{}\_v0.1\textbackslash{}\_instruct} (Instruct, N=1000)}

\subparagraph{【サマリー】 3×3 相関 \& 誤差マトリクス}
\begin{center}
\small
\begin{tabular}{lcccc}
\toprule
\textbf{タスク \ 感情次元} & \textbf{Valence ($V$)} & \textbf{Arousal ($A$)} & \textbf{Dominance ($D$)} & \textbf{完全中立率 $(5,5,5)$} \\
\midrule
\textbf{① Writer-State Estimation (W)} & \textbf{$r=0.486$} (MAE 0.43) & \textbf{$r=0.215$} (MAE 0.89) & \textbf{$r=0.055$} (MAE 0.65) & \texttt{0.6\textbackslash{}\%} \\
\textbf{② Reader-Response Prediction (R)} & \textbf{$r=0.496$} (MAE 0.44) & \textbf{$r=0.342$} (MAE 1.03) & \textbf{$r=0.002$} (MAE 0.79) & \texttt{0.0\textbackslash{}\%} \\
\textbf{③ Self-Report (S)} & \textbf{$r=0.551$} (MAE 0.43) & \textbf{$r=0.355$} (MAE 0.90) & \textbf{$r=0.062$} (MAE 0.60) & \texttt{0.1\textbackslash{}\%} \\
\bottomrule
\end{tabular}
\end{center}


\begin{itemize}
\item \textbf{内部結合度 (Internal Coupling $\text{Corr}(R, S)$)}: Valence $r=0.958$, Arousal $r=0.955$, Dominance $r=0.927
\end{itemize}

\subparagraph{【9個の詳細表】 タスク別・次元別詳細メトリクス}

\subparagraph{① Writer-State Estimation (W)}
\textit{評価基準: human\_writer VAD}

\textbf{表 1: ① Writer-State Estimation (W) — Valence}
\begin{center}
\small
\begin{tabular}{lcc}
\toprule
\textbf{評価軸 / 指標 (Metric)} & \textbf{値 (Value)} & \textbf{学術的意味・備考} \\
\midrule
\textbf{[Alignment] 連続相関 ($r$)} & \textbf{0.4863} & $p = <.001 $ (刺激間パターンの追従度) \\
\textbf{[Alignment] 順位相関 ($\rho$)} & 0.5181 & $p = <.001 $ (単調増加性) \\
\textbf{[Alignment] Greedy相関 ($r_{greedy}$)} & 0.3013 & 最尤トークン出力ベースの一致度 \\
\textbf{[Calibration] 平均絶対誤差 (MAE)} & \textbf{0.4290} & 人間正解値との絶対スケール距離 \\
\textbf{[Calibration] 二乗平均平方根誤差 (RMSE)} & 0.5277 & 大きな乖離に対する感度 \\
\textbf{[分布] モデル予測期待値平均 $\pm$ SD} & 2.734 $\pm$ 0.529 & $E[V]$ の連続分布 \\
\textbf{[分布] 人間アノテーション平均 $\pm$ SD} & 2.976 $\pm$ 0.335 & 人間正解値の分布 \\
\textbf{[Collapse] 単一次元 Greedy=5 固定率} & \texttt{31.30\textbackslash{}\%} & 当該次元の最頻出力が中立に退化した割合 \\
\bottomrule
\end{tabular}
\end{center}


\textbf{表 2: ① Writer-State Estimation (W) — Arousal}
\begin{center}
\small
\begin{tabular}{lcc}
\toprule
\textbf{評価軸 / 指標 (Metric)} & \textbf{値 (Value)} & \textbf{学術的意味・備考} \\
\midrule
\textbf{[Alignment] 連続相関 ($r$)} & \textbf{0.2151} & $p = <.001 $ (刺激間パターンの追従度) \\
\textbf{[Alignment] 順位相関 ($\rho$)} & 0.1818 & $p = <.001 $ (単調増加性) \\
\textbf{[Alignment] Greedy相関 ($r_{greedy}$)} & 0.1409 & 最尤トークン出力ベースの一致度 \\
\textbf{[Calibration] 平均絶対誤差 (MAE)} & \textbf{0.8902} & 人間正解値との絶対スケール距離 \\
\textbf{[Calibration] 二乗平均平方根誤差 (RMSE)} & 0.9955 & 大きな乖離に対する感度 \\
\textbf{[分布] モデル予測期待値平均 $\pm$ SD} & 2.165 $\pm$ 0.485 & $E[A]$ の連続分布 \\
\textbf{[分布] 人間アノテーション平均 $\pm$ SD} & 3.017 $\pm$ 0.305 & 人間正解値の分布 \\
\textbf{[Collapse] 単一次元 Greedy=5 固定率} & \texttt{10.70\textbackslash{}\%} & 当該次元の最頻出力が中立に退化した割合 \\
\bottomrule
\end{tabular}
\end{center}


\textbf{表 3: ① Writer-State Estimation (W) — Dominance}
\begin{center}
\small
\begin{tabular}{lcc}
\toprule
\textbf{評価軸 / 指標 (Metric)} & \textbf{値 (Value)} & \textbf{学術的意味・備考} \\
\midrule
\textbf{[Alignment] 連続相関 ($r$)} & \textbf{0.0551} & $p = 0.082 $ (刺激間パターンの追従度) \\
\textbf{[Alignment] 順位相関 ($\rho$)} & 0.0492 & $p = 0.120 $ (単調増加性) \\
\textbf{[Alignment] Greedy相関 ($r_{greedy}$)} & 0.0489 & 最尤トークン出力ベースの一致度 \\
\textbf{[Calibration] 平均絶対誤差 (MAE)} & \textbf{0.6466} & 人間正解値との絶対スケール距離 \\
\textbf{[Calibration] 二乗平均平方根誤差 (RMSE)} & 0.7295 & 大きな乖離に対する感度 \\
\textbf{[分布] モデル予測期待値平均 $\pm$ SD} & 2.463 $\pm$ 0.287 & $E[D]$ の連続分布 \\
\textbf{[分布] 人間アノテーション平均 $\pm$ SD} & 3.092 $\pm$ 0.251 & 人間正解値の分布 \\
\textbf{[Collapse] 単一次元 Greedy=5 固定率} & \texttt{15.40\textbackslash{}\%} & 当該次元の最頻出力が中立に退化した割合 \\
\bottomrule
\end{tabular}
\end{center}


\subparagraph{② Reader-Response Prediction (R)}
\textit{評価基準: human\_reader VAD}

\textbf{表 4: ② Reader-Response Prediction (R) — Valence}
\begin{center}
\small
\begin{tabular}{lcc}
\toprule
\textbf{評価軸 / 指標 (Metric)} & \textbf{値 (Value)} & \textbf{学術的意味・備考} \\
\midrule
\textbf{[Alignment] 連続相関 ($r$)} & \textbf{0.4956} & $p = <.001 $ (刺激間パターンの追従度) \\
\textbf{[Alignment] 順位相関 ($\rho$)} & 0.5328 & $p = <.001 $ (単調増加性) \\
\textbf{[Alignment] Greedy相関 ($r_{greedy}$)} & 0.3130 & 最尤トークン出力ベースの一致度 \\
\textbf{[Calibration] 平均絶対誤差 (MAE)} & \textbf{0.4421} & 人間正解値との絶対スケール距離 \\
\textbf{[Calibration] 二乗平均平方根誤差 (RMSE)} & 0.5492 & 大きな乖離に対する感度 \\
\textbf{[分布] モデル予測期待値平均 $\pm$ SD} & 2.716 $\pm$ 0.519 & $E[V]$ の連続分布 \\
\textbf{[分布] 人間アノテーション平均 $\pm$ SD} & 2.978 $\pm$ 0.428 & 人間正解値の分布 \\
\textbf{[Collapse] 単一次元 Greedy=5 固定率} & \texttt{41.00\textbackslash{}\%} & 当該次元の最頻出力が中立に退化した割合 \\
\bottomrule
\end{tabular}
\end{center}


\textbf{表 5: ② Reader-Response Prediction (R) — Arousal}
\begin{center}
\small
\begin{tabular}{lcc}
\toprule
\textbf{評価軸 / 指標 (Metric)} & \textbf{値 (Value)} & \textbf{学術的意味・備考} \\
\midrule
\textbf{[Alignment] 連続相関 ($r$)} & \textbf{0.3423} & $p = <.001 $ (刺激間パターンの追従度) \\
\textbf{[Alignment] 順位相関 ($\rho$)} & 0.3026 & $p = <.001 $ (単調増加性) \\
\textbf{[Alignment] Greedy相関 ($r_{greedy}$)} & 0.3083 & 最尤トークン出力ベースの一致度 \\
\textbf{[Calibration] 平均絶対誤差 (MAE)} & \textbf{1.0267} & 人間正解値との絶対スケール距離 \\
\textbf{[Calibration] 二乗平均平方根誤差 (RMSE)} & 1.1225 & 大きな乖離に対する感度 \\
\textbf{[分布] モデル予測期待値平均 $\pm$ SD} & 2.038 $\pm$ 0.492 & $E[A]$ の連続分布 \\
\textbf{[分布] 人間アノテーション平均 $\pm$ SD} & 3.050 $\pm$ 0.313 & 人間正解値の分布 \\
\textbf{[Collapse] 単一次元 Greedy=5 固定率} & \texttt{8.10\textbackslash{}\%} & 当該次元の最頻出力が中立に退化した割合 \\
\bottomrule
\end{tabular}
\end{center}


\textbf{表 6: ② Reader-Response Prediction (R) — Dominance}
\begin{center}
\small
\begin{tabular}{lcc}
\toprule
\textbf{評価軸 / 指標 (Metric)} & \textbf{値 (Value)} & \textbf{学術的意味・備考} \\
\midrule
\textbf{[Alignment] 連続相関 ($r$)} & \textbf{0.0018} & $p = 0.954 $ (刺激間パターンの追従度) \\
\textbf{[Alignment] 順位相関 ($\rho$)} & -0.0133 & $p = 0.676 $ (単調増加性) \\
\textbf{[Alignment] Greedy相関 ($r_{greedy}$)} & -0.0036 & 最尤トークン出力ベースの一致度 \\
\textbf{[Calibration] 平均絶対誤差 (MAE)} & \textbf{0.7877} & 人間正解値との絶対スケール距離 \\
\textbf{[Calibration] 二乗平均平方根誤差 (RMSE)} & 0.8735 & 大きな乖離に対する感度 \\
\textbf{[分布] モデル予測期待値平均 $\pm$ SD} & 2.272 $\pm$ 0.296 & $E[D]$ の連続分布 \\
\textbf{[分布] 人間アノテーション平均 $\pm$ SD} & 3.043 $\pm$ 0.287 & 人間正解値の分布 \\
\textbf{[Collapse] 単一次元 Greedy=5 固定率} & \texttt{9.30\textbackslash{}\%} & 当該次元の最頻出力が中立に退化した割合 \\
\bottomrule
\end{tabular}
\end{center}


\subparagraph{③ Self-Report (S)}
\textit{評価基準: human\_reader VAD}

\textbf{表 7: ③ Self-Report (S) — Valence}
\begin{center}
\small
\begin{tabular}{lcc}
\toprule
\textbf{評価軸 / 指標 (Metric)} & \textbf{値 (Value)} & \textbf{学術的意味・備考} \\
\midrule
\textbf{[Alignment] 連続相関 ($r$)} & \textbf{0.5507} & $p = <.001 $ (刺激間パターンの追従度) \\
\textbf{[Alignment] 順位相関 ($\rho$)} & 0.5911 & $p = <.001 $ (単調増加性) \\
\textbf{[Alignment] Greedy相関 ($r_{greedy}$)} & 0.4024 & 最尤トークン出力ベースの一致度 \\
\textbf{[Calibration] 平均絶対誤差 (MAE)} & \textbf{0.4261} & 人間正解値との絶対スケール距離 \\
\textbf{[Calibration] 二乗平均平方根誤差 (RMSE)} & 0.5325 & 大きな乖離に対する感度 \\
\textbf{[分布] モデル予測期待値平均 $\pm$ SD} & 2.808 $\pm$ 0.592 & $E[V]$ の連続分布 \\
\textbf{[分布] 人間アノテーション平均 $\pm$ SD} & 2.978 $\pm$ 0.428 & 人間正解値の分布 \\
\textbf{[Collapse] 単一次元 Greedy=5 固定率} & \texttt{28.10\textbackslash{}\%} & 当該次元の最頻出力が中立に退化した割合 \\
\bottomrule
\end{tabular}
\end{center}


\textbf{表 8: ③ Self-Report (S) — Arousal}
\begin{center}
\small
\begin{tabular}{lcc}
\toprule
\textbf{評価軸 / 指標 (Metric)} & \textbf{値 (Value)} & \textbf{学術的意味・備考} \\
\midrule
\textbf{[Alignment] 連続相関 ($r$)} & \textbf{0.3546} & $p = <.001 $ (刺激間パターンの追従度) \\
\textbf{[Alignment] 順位相関 ($\rho$)} & 0.3136 & $p = <.001 $ (単調増加性) \\
\textbf{[Alignment] Greedy相関 ($r_{greedy}$)} & 0.2972 & 最尤トークン出力ベースの一致度 \\
\textbf{[Calibration] 平均絶対誤差 (MAE)} & \textbf{0.9028} & 人間正解値との絶対スケール距離 \\
\textbf{[Calibration] 二乗平均平方根誤差 (RMSE)} & 1.0019 & 大きな乖離に対する感度 \\
\textbf{[分布] モデル予測期待値平均 $\pm$ SD} & 2.165 $\pm$ 0.478 & $E[A]$ の連続分布 \\
\textbf{[分布] 人間アノテーション平均 $\pm$ SD} & 3.050 $\pm$ 0.313 & 人間正解値の分布 \\
\textbf{[Collapse] 単一次元 Greedy=5 固定率} & \texttt{11.60\textbackslash{}\%} & 当該次元の最頻出力が中立に退化した割合 \\
\bottomrule
\end{tabular}
\end{center}


\textbf{表 9: ③ Self-Report (S) — Dominance}
\begin{center}
\small
\begin{tabular}{lcc}
\toprule
\textbf{評価軸 / 指標 (Metric)} & \textbf{値 (Value)} & \textbf{学術的意味・備考} \\
\midrule
\textbf{[Alignment] 連続相関 ($r$)} & \textbf{0.0619} & $p = 0.050 $ (刺激間パターンの追従度) \\
\textbf{[Alignment] 順位相関 ($\rho$)} & 0.0468 & $p = 0.139 $ (単調増加性) \\
\textbf{[Alignment] Greedy相関 ($r_{greedy}$)} & 0.0749 & 最尤トークン出力ベースの一致度 \\
\textbf{[Calibration] 平均絶対誤差 (MAE)} & \textbf{0.6004} & 人間正解値との絶対スケール距離 \\
\textbf{[Calibration] 二乗平均平方根誤差 (RMSE)} & 0.6901 & 大きな乖離に対する感度 \\
\textbf{[分布] モデル予測期待値平均 $\pm$ SD} & 2.471 $\pm$ 0.277 & $E[D]$ の連続分布 \\
\textbf{[分布] 人間アノテーション平均 $\pm$ SD} & 3.043 $\pm$ 0.287 & 人間正解値の分布 \\
\textbf{[Collapse] 単一次元 Greedy=5 固定率} & \texttt{21.70\textbackslash{}\%} & 当該次元の最頻出力が中立に退化した割合 \\
\bottomrule
\end{tabular}
\end{center}


\vspace{0.5em}\hrule\vspace{0.5em}

\paragraph{■ モデル: \texttt{qwen2.5\textbackslash{}\_1.5b\textbackslash{}\_base} (Base, N=1000)}

\subparagraph{【サマリー】 3×3 相関 \& 誤差マトリクス}
\begin{center}
\small
\begin{tabular}{lcccc}
\toprule
\textbf{タスク \ 感情次元} & \textbf{Valence ($V$)} & \textbf{Arousal ($A$)} & \textbf{Dominance ($D$)} & \textbf{完全中立率 $(5,5,5)$} \\
\midrule
\textbf{① Writer-State Estimation (W)} & \textbf{$r=0.542$} (MAE 0.33) & \textbf{$r=0.159$} (MAE 0.33) & \textbf{$r=0.098$} (MAE 0.30) & \texttt{0.1\textbackslash{}\%} \\
\textbf{② Reader-Response Prediction (R)} & \textbf{$r=0.573$} (MAE 0.35) & \textbf{$r=0.248$} (MAE 0.31) & \textbf{$r=0.225$} (MAE 0.33) & \texttt{0.0\textbackslash{}\%} \\
\textbf{③ Self-Report (S)} & \textbf{$r=0.588$} (MAE 0.30) & \textbf{$r=0.191$} (MAE 0.34) & \textbf{$r=0.224$} (MAE 0.37) & \texttt{0.0\textbackslash{}\%} \\
\bottomrule
\end{tabular}
\end{center}


\begin{itemize}
\item \textbf{内部結合度 (Internal Coupling $\text{Corr}(R, S)$)}: Valence $r=0.962$, Arousal $r=0.942$, Dominance $r=0.946
\end{itemize}

\subparagraph{【9個の詳細表】 タスク別・次元別詳細メトリクス}

\subparagraph{① Writer-State Estimation (W)}
\textit{評価基準: human\_writer VAD}

\textbf{表 1: ① Writer-State Estimation (W) — Valence}
\begin{center}
\small
\begin{tabular}{lcc}
\toprule
\textbf{評価軸 / 指標 (Metric)} & \textbf{値 (Value)} & \textbf{学術的意味・備考} \\
\midrule
\textbf{[Alignment] 連続相関 ($r$)} & \textbf{0.5423} & $p = <.001 $ (刺激間パターンの追従度) \\
\textbf{[Alignment] 順位相関 ($\rho$)} & 0.5746 & $p = <.001 $ (単調増加性) \\
\textbf{[Alignment] Greedy相関 ($r_{greedy}$)} & 0.3751 & 最尤トークン出力ベースの一致度 \\
\textbf{[Calibration] 平均絶対誤差 (MAE)} & \textbf{0.3273} & 人間正解値との絶対スケール距離 \\
\textbf{[Calibration] 二乗平均平方根誤差 (RMSE)} & 0.4272 & 大きな乖離に対する感度 \\
\textbf{[分布] モデル予測期待値平均 $\pm$ SD} & 2.812 $\pm$ 0.458 & $E[V]$ の連続分布 \\
\textbf{[分布] 人間アノテーション平均 $\pm$ SD} & 2.976 $\pm$ 0.335 & 人間正解値の分布 \\
\textbf{[Collapse] 単一次元 Greedy=5 固定率} & \texttt{0.10\textbackslash{}\%} & 当該次元の最頻出力が中立に退化した割合 \\
\bottomrule
\end{tabular}
\end{center}


\textbf{表 2: ① Writer-State Estimation (W) — Arousal}
\begin{center}
\small
\begin{tabular}{lcc}
\toprule
\textbf{評価軸 / 指標 (Metric)} & \textbf{値 (Value)} & \textbf{学術的意味・備考} \\
\midrule
\textbf{[Alignment] 連続相関 ($r$)} & \textbf{0.1586} & $p = <.001 $ (刺激間パターンの追従度) \\
\textbf{[Alignment] 順位相関 ($\rho$)} & 0.1451 & $p = <.001 $ (単調増加性) \\
\textbf{[Alignment] Greedy相関 ($r_{greedy}$)} & 0.1695 & 最尤トークン出力ベースの一致度 \\
\textbf{[Calibration] 平均絶対誤差 (MAE)} & \textbf{0.3272} & 人間正解値との絶対スケール距離 \\
\textbf{[Calibration] 二乗平均平方根誤差 (RMSE)} & 0.4141 & 大きな乖離に対する感度 \\
\textbf{[分布] モデル予測期待値平均 $\pm$ SD} & 3.109 $\pm$ 0.318 & $E[A]$ の連続分布 \\
\textbf{[分布] 人間アノテーション平均 $\pm$ SD} & 3.017 $\pm$ 0.305 & 人間正解値の分布 \\
\textbf{[Collapse] 単一次元 Greedy=5 固定率} & \texttt{0.10\textbackslash{}\%} & 当該次元の最頻出力が中立に退化した割合 \\
\bottomrule
\end{tabular}
\end{center}


\textbf{表 3: ① Writer-State Estimation (W) — Dominance}
\begin{center}
\small
\begin{tabular}{lcc}
\toprule
\textbf{評価軸 / 指標 (Metric)} & \textbf{値 (Value)} & \textbf{学術的意味・備考} \\
\midrule
\textbf{[Alignment] 連続相関 ($r$)} & \textbf{0.0982} & $p = 0.002 $ (刺激間パターンの追従度) \\
\textbf{[Alignment] 順位相関 ($\rho$)} & 0.1069 & $p = <.001 $ (単調増加性) \\
\textbf{[Alignment] Greedy相関 ($r_{greedy}$)} & 0.1290 & 最尤トークン出力ベースの一致度 \\
\textbf{[Calibration] 平均絶対誤差 (MAE)} & \textbf{0.3031} & 人間正解値との絶対スケール距離 \\
\textbf{[Calibration] 二乗平均平方根誤差 (RMSE)} & 0.3757 & 大きな乖離に対する感度 \\
\textbf{[分布] モデル予測期待値平均 $\pm$ SD} & 3.208 $\pm$ 0.280 & $E[D]$ の連続分布 \\
\textbf{[分布] 人間アノテーション平均 $\pm$ SD} & 3.092 $\pm$ 0.251 & 人間正解値の分布 \\
\textbf{[Collapse] 単一次元 Greedy=5 固定率} & \texttt{0.10\textbackslash{}\%} & 当該次元の最頻出力が中立に退化した割合 \\
\bottomrule
\end{tabular}
\end{center}


\subparagraph{② Reader-Response Prediction (R)}
\textit{評価基準: human\_reader VAD}

\textbf{表 4: ② Reader-Response Prediction (R) — Valence}
\begin{center}
\small
\begin{tabular}{lcc}
\toprule
\textbf{評価軸 / 指標 (Metric)} & \textbf{値 (Value)} & \textbf{学術的意味・備考} \\
\midrule
\textbf{[Alignment] 連続相関 ($r$)} & \textbf{0.5727} & $p = <.001 $ (刺激間パターンの追従度) \\
\textbf{[Alignment] 順位相関 ($\rho$)} & 0.5906 & $p = <.001 $ (単調増加性) \\
\textbf{[Alignment] Greedy相関 ($r_{greedy}$)} & 0.3497 & 最尤トークン出力ベースの一致度 \\
\textbf{[Calibration] 平均絶対誤差 (MAE)} & \textbf{0.3507} & 人間正解値との絶対スケール距離 \\
\textbf{[Calibration] 二乗平均平方根誤差 (RMSE)} & 0.4559 & 大きな乖離に対する感度 \\
\textbf{[分布] モデル予測期待値平均 $\pm$ SD} & 2.823 $\pm$ 0.491 & $E[V]$ の連続分布 \\
\textbf{[分布] 人間アノテーション平均 $\pm$ SD} & 2.978 $\pm$ 0.428 & 人間正解値の分布 \\
\textbf{[Collapse] 単一次元 Greedy=5 固定率} & \texttt{0.00\textbackslash{}\%} & 当該次元の最頻出力が中立に退化した割合 \\
\bottomrule
\end{tabular}
\end{center}


\textbf{表 5: ② Reader-Response Prediction (R) — Arousal}
\begin{center}
\small
\begin{tabular}{lcc}
\toprule
\textbf{評価軸 / 指標 (Metric)} & \textbf{値 (Value)} & \textbf{学術的意味・備考} \\
\midrule
\textbf{[Alignment] 連続相関 ($r$)} & \textbf{0.2476} & $p = <.001 $ (刺激間パターンの追従度) \\
\textbf{[Alignment] 順位相関 ($\rho$)} & 0.2088 & $p = <.001 $ (単調増加性) \\
\textbf{[Alignment] Greedy相関 ($r_{greedy}$)} & 0.1344 & 最尤トークン出力ベースの一致度 \\
\textbf{[Calibration] 平均絶対誤差 (MAE)} & \textbf{0.3112} & 人間正解値との絶対スケール距離 \\
\textbf{[Calibration] 二乗平均平方根誤差 (RMSE)} & 0.3963 & 大きな乖離に対する感度 \\
\textbf{[分布] モデル予測期待値平均 $\pm$ SD} & 3.101 $\pm$ 0.327 & $E[A]$ の連続分布 \\
\textbf{[分布] 人間アノテーション平均 $\pm$ SD} & 3.050 $\pm$ 0.313 & 人間正解値の分布 \\
\textbf{[Collapse] 単一次元 Greedy=5 固定率} & \texttt{0.70\textbackslash{}\%} & 当該次元の最頻出力が中立に退化した割合 \\
\bottomrule
\end{tabular}
\end{center}


\textbf{表 6: ② Reader-Response Prediction (R) — Dominance}
\begin{center}
\small
\begin{tabular}{lcc}
\toprule
\textbf{評価軸 / 指標 (Metric)} & \textbf{値 (Value)} & \textbf{学術的意味・備考} \\
\midrule
\textbf{[Alignment] 連続相関 ($r$)} & \textbf{0.2255} & $p = <.001 $ (刺激間パターンの追従度) \\
\textbf{[Alignment] 順位相関 ($\rho$)} & 0.2177 & $p = <.001 $ (単調増加性) \\
\textbf{[Alignment] Greedy相関 ($r_{greedy}$)} & 0.1602 & 最尤トークン出力ベースの一致度 \\
\textbf{[Calibration] 平均絶対誤差 (MAE)} & \textbf{0.3283} & 人間正解値との絶対スケール距離 \\
\textbf{[Calibration] 二乗平均平方根誤差 (RMSE)} & 0.4013 & 大きな乖離に対する感度 \\
\textbf{[分布] モデル予測期待値平均 $\pm$ SD} & 3.230 $\pm$ 0.284 & $E[D]$ の連続分布 \\
\textbf{[分布] 人間アノテーション平均 $\pm$ SD} & 3.043 $\pm$ 0.287 & 人間正解値の分布 \\
\textbf{[Collapse] 単一次元 Greedy=5 固定率} & \texttt{0.00\textbackslash{}\%} & 当該次元の最頻出力が中立に退化した割合 \\
\bottomrule
\end{tabular}
\end{center}


\subparagraph{③ Self-Report (S)}
\textit{評価基準: human\_reader VAD}

\textbf{表 7: ③ Self-Report (S) — Valence}
\begin{center}
\small
\begin{tabular}{lcc}
\toprule
\textbf{評価軸 / 指標 (Metric)} & \textbf{値 (Value)} & \textbf{学術的意味・備考} \\
\midrule
\textbf{[Alignment] 連続相関 ($r$)} & \textbf{0.5883} & $p = <.001 $ (刺激間パターンの追従度) \\
\textbf{[Alignment] 順位相関 ($\rho$)} & 0.6102 & $p = <.001 $ (単調増加性) \\
\textbf{[Alignment] Greedy相関 ($r_{greedy}$)} & 0.3624 & 最尤トークン出力ベースの一致度 \\
\textbf{[Calibration] 平均絶対誤差 (MAE)} & \textbf{0.3011} & 人間正解値との絶対スケール距離 \\
\textbf{[Calibration] 二乗平均平方根誤差 (RMSE)} & 0.4001 & 大きな乖離に対する感度 \\
\textbf{[分布] モデル予測期待値平均 $\pm$ SD} & 2.893 $\pm$ 0.433 & $E[V]$ の連続分布 \\
\textbf{[分布] 人間アノテーション平均 $\pm$ SD} & 2.978 $\pm$ 0.428 & 人間正解値の分布 \\
\textbf{[Collapse] 単一次元 Greedy=5 固定率} & \texttt{0.10\textbackslash{}\%} & 当該次元の最頻出力が中立に退化した割合 \\
\bottomrule
\end{tabular}
\end{center}


\textbf{表 8: ③ Self-Report (S) — Arousal}
\begin{center}
\small
\begin{tabular}{lcc}
\toprule
\textbf{評価軸 / 指標 (Metric)} & \textbf{値 (Value)} & \textbf{学術的意味・備考} \\
\midrule
\textbf{[Alignment] 連続相関 ($r$)} & \textbf{0.1913} & $p = <.001 $ (刺激間パターンの追従度) \\
\textbf{[Alignment] 順位相関 ($\rho$)} & 0.1496 & $p = <.001 $ (単調増加性) \\
\textbf{[Alignment] Greedy相関 ($r_{greedy}$)} & 0.1346 & 最尤トークン出力ベースの一致度 \\
\textbf{[Calibration] 平均絶対誤差 (MAE)} & \textbf{0.3427} & 人間正解値との絶対スケール距離 \\
\textbf{[Calibration] 二乗平均平方根誤差 (RMSE)} & 0.4287 & 大きな乖離に対する感度 \\
\textbf{[分布] モデル予測期待値平均 $\pm$ SD} & 3.212 $\pm$ 0.311 & $E[A]$ の連続分布 \\
\textbf{[分布] 人間アノテーション平均 $\pm$ SD} & 3.050 $\pm$ 0.313 & 人間正解値の分布 \\
\textbf{[Collapse] 単一次元 Greedy=5 固定率} & \texttt{0.10\textbackslash{}\%} & 当該次元の最頻出力が中立に退化した割合 \\
\bottomrule
\end{tabular}
\end{center}


\textbf{表 9: ③ Self-Report (S) — Dominance}
\begin{center}
\small
\begin{tabular}{lcc}
\toprule
\textbf{評価軸 / 指標 (Metric)} & \textbf{値 (Value)} & \textbf{学術的意味・備考} \\
\midrule
\textbf{[Alignment] 連続相関 ($r$)} & \textbf{0.2239} & $p = <.001 $ (刺激間パターンの追従度) \\
\textbf{[Alignment] 順位相関 ($\rho$)} & 0.2144 & $p = <.001 $ (単調増加性) \\
\textbf{[Alignment] Greedy相関 ($r_{greedy}$)} & 0.1496 & 最尤トークン出力ベースの一致度 \\
\textbf{[Calibration] 平均絶対誤差 (MAE)} & \textbf{0.3671} & 人間正解値との絶対スケール距離 \\
\textbf{[Calibration] 二乗平均平方根誤差 (RMSE)} & 0.4408 & 大きな乖離に対する感度 \\
\textbf{[分布] モデル予測期待値平均 $\pm$ SD} & 3.304 $\pm$ 0.283 & $E[D]$ の連続分布 \\
\textbf{[分布] 人間アノテーション平均 $\pm$ SD} & 3.043 $\pm$ 0.287 & 人間正解値の分布 \\
\textbf{[Collapse] 単一次元 Greedy=5 固定率} & \texttt{0.00\textbackslash{}\%} & 当該次元の最頻出力が中立に退化した割合 \\
\bottomrule
\end{tabular}
\end{center}


\vspace{0.5em}\hrule\vspace{0.5em}

\paragraph{■ モデル: \texttt{qwen2.5\textbackslash{}\_1.5b\textbackslash{}\_instruct} (Instruct, N=1000)}

\subparagraph{【サマリー】 3×3 相関 \& 誤差マトリクス}
\begin{center}
\small
\begin{tabular}{lcccc}
\toprule
\textbf{タスク \ 感情次元} & \textbf{Valence ($V$)} & \textbf{Arousal ($A$)} & \textbf{Dominance ($D$)} & \textbf{完全中立率 $(5,5,5)$} \\
\midrule
\textbf{① Writer-State Estimation (W)} & \textbf{$r=0.553$} (MAE 0.57) & \textbf{$r=0.090$} (MAE 0.66) & \textbf{$r=0.036$} (MAE 0.52) & \texttt{0.0\textbackslash{}\%} \\
\textbf{② Reader-Response Prediction (R)} & \textbf{$r=0.590$} (MAE 0.58) & \textbf{$r=0.188$} (MAE 0.70) & \textbf{$r=0.122$} (MAE 0.54) & \texttt{0.0\textbackslash{}\%} \\
\textbf{③ Self-Report (S)} & \textbf{$r=0.641$} (MAE 0.49) & \textbf{$r=0.125$} (MAE 0.51) & \textbf{$r=0.165$} (MAE 0.44) & \texttt{0.4\textbackslash{}\%} \\
\bottomrule
\end{tabular}
\end{center}


\begin{itemize}
\item \textbf{内部結合度 (Internal Coupling $\text{Corr}(R, S)$)}: Valence $r=0.944$, Arousal $r=0.872$, Dominance $r=0.894
\end{itemize}

\subparagraph{【9個の詳細表】 タスク別・次元別詳細メトリクス}

\subparagraph{① Writer-State Estimation (W)}
\textit{評価基準: human\_writer VAD}

\textbf{表 1: ① Writer-State Estimation (W) — Valence}
\begin{center}
\small
\begin{tabular}{lcc}
\toprule
\textbf{評価軸 / 指標 (Metric)} & \textbf{値 (Value)} & \textbf{学術的意味・備考} \\
\midrule
\textbf{[Alignment] 連続相関 ($r$)} & \textbf{0.5533} & $p = <.001 $ (刺激間パターンの追従度) \\
\textbf{[Alignment] 順位相関 ($\rho$)} & 0.5763 & $p = <.001 $ (単調増加性) \\
\textbf{[Alignment] Greedy相関 ($r_{greedy}$)} & 0.3474 & 最尤トークン出力ベースの一致度 \\
\textbf{[Calibration] 平均絶対誤差 (MAE)} & \textbf{0.5684} & 人間正解値との絶対スケール距離 \\
\textbf{[Calibration] 二乗平均平方根誤差 (RMSE)} & 0.6828 & 大きな乖離に対する感度 \\
\textbf{[分布] モデル予測期待値平均 $\pm$ SD} & 2.762 $\pm$ 0.771 & $E[V]$ の連続分布 \\
\textbf{[分布] 人間アノテーション平均 $\pm$ SD} & 2.976 $\pm$ 0.335 & 人間正解値の分布 \\
\textbf{[Collapse] 単一次元 Greedy=5 固定率} & \texttt{0.30\textbackslash{}\%} & 当該次元の最頻出力が中立に退化した割合 \\
\bottomrule
\end{tabular}
\end{center}


\textbf{表 2: ① Writer-State Estimation (W) — Arousal}
\begin{center}
\small
\begin{tabular}{lcc}
\toprule
\textbf{評価軸 / 指標 (Metric)} & \textbf{値 (Value)} & \textbf{学術的意味・備考} \\
\midrule
\textbf{[Alignment] 連続相関 ($r$)} & \textbf{0.0901} & $p = 0.004 $ (刺激間パターンの追従度) \\
\textbf{[Alignment] 順位相関 ($\rho$)} & 0.0637 & $p = 0.044 $ (単調増加性) \\
\textbf{[Alignment] Greedy相関 ($r_{greedy}$)} & 0.0289 & 最尤トークン出力ベースの一致度 \\
\textbf{[Calibration] 平均絶対誤差 (MAE)} & \textbf{0.6555} & 人間正解値との絶対スケール距離 \\
\textbf{[Calibration] 二乗平均平方根誤差 (RMSE)} & 0.7701 & 大きな乖離に対する感度 \\
\textbf{[分布] モデル予測期待値平均 $\pm$ SD} & 2.389 $\pm$ 0.354 & $E[A]$ の連続分布 \\
\textbf{[分布] 人間アノテーション平均 $\pm$ SD} & 3.017 $\pm$ 0.305 & 人間正解値の分布 \\
\textbf{[Collapse] 単一次元 Greedy=5 固定率} & \texttt{11.70\textbackslash{}\%} & 当該次元の最頻出力が中立に退化した割合 \\
\bottomrule
\end{tabular}
\end{center}


\textbf{表 3: ① Writer-State Estimation (W) — Dominance}
\begin{center}
\small
\begin{tabular}{lcc}
\toprule
\textbf{評価軸 / 指標 (Metric)} & \textbf{値 (Value)} & \textbf{学術的意味・備考} \\
\midrule
\textbf{[Alignment] 連続相関 ($r$)} & \textbf{0.0361} & $p = 0.255 $ (刺激間パターンの追従度) \\
\textbf{[Alignment] 順位相関 ($\rho$)} & 0.0578 & $p = 0.068 $ (単調増加性) \\
\textbf{[Alignment] Greedy相関 ($r_{greedy}$)} & 0.0322 & 最尤トークン出力ベースの一致度 \\
\textbf{[Calibration] 平均絶対誤差 (MAE)} & \textbf{0.5192} & 人間正解値との絶対スケール距離 \\
\textbf{[Calibration] 二乗平均平方根誤差 (RMSE)} & 0.6310 & 大きな乖離に対する感度 \\
\textbf{[分布] モデル予測期待値平均 $\pm$ SD} & 2.598 $\pm$ 0.311 & $E[D]$ の連続分布 \\
\textbf{[分布] 人間アノテーション平均 $\pm$ SD} & 3.092 $\pm$ 0.251 & 人間正解値の分布 \\
\textbf{[Collapse] 単一次元 Greedy=5 固定率} & \texttt{11.10\textbackslash{}\%} & 当該次元の最頻出力が中立に退化した割合 \\
\bottomrule
\end{tabular}
\end{center}


\subparagraph{② Reader-Response Prediction (R)}
\textit{評価基準: human\_reader VAD}

\textbf{表 4: ② Reader-Response Prediction (R) — Valence}
\begin{center}
\small
\begin{tabular}{lcc}
\toprule
\textbf{評価軸 / 指標 (Metric)} & \textbf{値 (Value)} & \textbf{学術的意味・備考} \\
\midrule
\textbf{[Alignment] 連続相関 ($r$)} & \textbf{0.5903} & $p = <.001 $ (刺激間パターンの追従度) \\
\textbf{[Alignment] 順位相関 ($\rho$)} & 0.6071 & $p = <.001 $ (単調増加性) \\
\textbf{[Alignment] Greedy相関 ($r_{greedy}$)} & 0.3895 & 最尤トークン出力ベースの一致度 \\
\textbf{[Calibration] 平均絶対誤差 (MAE)} & \textbf{0.5790} & 人間正解値との絶対スケール距離 \\
\textbf{[Calibration] 二乗平均平方根誤差 (RMSE)} & 0.7036 & 大きな乖離に対する感度 \\
\textbf{[分布] モデル予測期待値平均 $\pm$ SD} & 2.754 $\pm$ 0.823 & $E[V]$ の連続分布 \\
\textbf{[分布] 人間アノテーション平均 $\pm$ SD} & 2.978 $\pm$ 0.428 & 人間正解値の分布 \\
\textbf{[Collapse] 単一次元 Greedy=5 固定率} & \texttt{0.60\textbackslash{}\%} & 当該次元の最頻出力が中立に退化した割合 \\
\bottomrule
\end{tabular}
\end{center}


\textbf{表 5: ② Reader-Response Prediction (R) — Arousal}
\begin{center}
\small
\begin{tabular}{lcc}
\toprule
\textbf{評価軸 / 指標 (Metric)} & \textbf{値 (Value)} & \textbf{学術的意味・備考} \\
\midrule
\textbf{[Alignment] 連続相関 ($r$)} & \textbf{0.1876} & $p = <.001 $ (刺激間パターンの追従度) \\
\textbf{[Alignment] 順位相関 ($\rho$)} & 0.1378 & $p = <.001 $ (単調増加性) \\
\textbf{[Alignment] Greedy相関 ($r_{greedy}$)} & -0.0117 & 最尤トークン出力ベースの一致度 \\
\textbf{[Calibration] 平均絶対誤差 (MAE)} & \textbf{0.6958} & 人間正解値との絶対スケール距離 \\
\textbf{[Calibration] 二乗平均平方根誤差 (RMSE)} & 0.8091 & 大きな乖離に対する感度 \\
\textbf{[分布] モデル予測期待値平均 $\pm$ SD} & 2.381 $\pm$ 0.393 & $E[A]$ の連続分布 \\
\textbf{[分布] 人間アノテーション平均 $\pm$ SD} & 3.050 $\pm$ 0.313 & 人間正解値の分布 \\
\textbf{[Collapse] 単一次元 Greedy=5 固定率} & \texttt{13.40\textbackslash{}\%} & 当該次元の最頻出力が中立に退化した割合 \\
\bottomrule
\end{tabular}
\end{center}


\textbf{表 6: ② Reader-Response Prediction (R) — Dominance}
\begin{center}
\small
\begin{tabular}{lcc}
\toprule
\textbf{評価軸 / 指標 (Metric)} & \textbf{値 (Value)} & \textbf{学術的意味・備考} \\
\midrule
\textbf{[Alignment] 連続相関 ($r$)} & \textbf{0.1224} & $p = <.001 $ (刺激間パターンの追従度) \\
\textbf{[Alignment] 順位相関 ($\rho$)} & 0.1305 & $p = <.001 $ (単調増加性) \\
\textbf{[Alignment] Greedy相関 ($r_{greedy}$)} & 0.1252 & 最尤トークン出力ベースの一致度 \\
\textbf{[Calibration] 平均絶対誤差 (MAE)} & \textbf{0.5432} & 人間正解値との絶対スケール距離 \\
\textbf{[Calibration] 二乗平均平方根誤差 (RMSE)} & 0.6603 & 大きな乖離に対する感度 \\
\textbf{[分布] モデル予測期待値平均 $\pm$ SD} & 2.533 $\pm$ 0.344 & $E[D]$ の連続分布 \\
\textbf{[分布] 人間アノテーション平均 $\pm$ SD} & 3.043 $\pm$ 0.287 & 人間正解値の分布 \\
\textbf{[Collapse] 単一次元 Greedy=5 固定率} & \texttt{15.00\textbackslash{}\%} & 当該次元の最頻出力が中立に退化した割合 \\
\bottomrule
\end{tabular}
\end{center}


\subparagraph{③ Self-Report (S)}
\textit{評価基準: human\_reader VAD}

\textbf{表 7: ③ Self-Report (S) — Valence}
\begin{center}
\small
\begin{tabular}{lcc}
\toprule
\textbf{評価軸 / 指標 (Metric)} & \textbf{値 (Value)} & \textbf{学術的意味・備考} \\
\midrule
\textbf{[Alignment] 連続相関 ($r$)} & \textbf{0.6411} & $p = <.001 $ (刺激間パターンの追従度) \\
\textbf{[Alignment] 順位相関 ($\rho$)} & 0.6498 & $p = <.001 $ (単調増加性) \\
\textbf{[Alignment] Greedy相関 ($r_{greedy}$)} & 0.4872 & 最尤トークン出力ベースの一致度 \\
\textbf{[Calibration] 平均絶対誤差 (MAE)} & \textbf{0.4912} & 人間正解値との絶対スケール距離 \\
\textbf{[Calibration] 二乗平均平方根誤差 (RMSE)} & 0.5983 & 大きな乖離に対する感度 \\
\textbf{[分布] モデル予測期待値平均 $\pm$ SD} & 2.990 $\pm$ 0.774 & $E[V]$ の連続分布 \\
\textbf{[分布] 人間アノテーション平均 $\pm$ SD} & 2.978 $\pm$ 0.428 & 人間正解値の分布 \\
\textbf{[Collapse] 単一次元 Greedy=5 固定率} & \texttt{2.30\textbackslash{}\%} & 当該次元の最頻出力が中立に退化した割合 \\
\bottomrule
\end{tabular}
\end{center}


\textbf{表 8: ③ Self-Report (S) — Arousal}
\begin{center}
\small
\begin{tabular}{lcc}
\toprule
\textbf{評価軸 / 指標 (Metric)} & \textbf{値 (Value)} & \textbf{学術的意味・備考} \\
\midrule
\textbf{[Alignment] 連続相関 ($r$)} & \textbf{0.1252} & $p = <.001 $ (刺激間パターンの追従度) \\
\textbf{[Alignment] 順位相関 ($\rho$)} & 0.0718 & $p = 0.023 $ (単調増加性) \\
\textbf{[Alignment] Greedy相関 ($r_{greedy}$)} & -0.0598 & 最尤トークン出力ベースの一致度 \\
\textbf{[Calibration] 平均絶対誤差 (MAE)} & \textbf{0.5077} & 人間正解値との絶対スケール距離 \\
\textbf{[Calibration] 二乗平均平方根誤差 (RMSE)} & 0.6302 & 大きな乖離に対する感度 \\
\textbf{[分布] モデル予測期待値平均 $\pm$ SD} & 2.625 $\pm$ 0.385 & $E[A]$ の連続分布 \\
\textbf{[分布] 人間アノテーション平均 $\pm$ SD} & 3.050 $\pm$ 0.313 & 人間正解値の分布 \\
\textbf{[Collapse] 単一次元 Greedy=5 固定率} & \texttt{34.00\textbackslash{}\%} & 当該次元の最頻出力が中立に退化した割合 \\
\bottomrule
\end{tabular}
\end{center}


\textbf{表 9: ③ Self-Report (S) — Dominance}
\begin{center}
\small
\begin{tabular}{lcc}
\toprule
\textbf{評価軸 / 指標 (Metric)} & \textbf{値 (Value)} & \textbf{学術的意味・備考} \\
\midrule
\textbf{[Alignment] 連続相関 ($r$)} & \textbf{0.1651} & $p = <.001 $ (刺激間パターンの追従度) \\
\textbf{[Alignment] 順位相関 ($\rho$)} & 0.1678 & $p = <.001 $ (単調増加性) \\
\textbf{[Alignment] Greedy相関 ($r_{greedy}$)} & 0.1536 & 最尤トークン出力ベースの一致度 \\
\textbf{[Calibration] 平均絶対誤差 (MAE)} & \textbf{0.4440} & 人間正解値との絶対スケール距離 \\
\textbf{[Calibration] 二乗平均平方根誤差 (RMSE)} & 0.5573 & 大きな乖離に対する感度 \\
\textbf{[分布] モデル予測期待値平均 $\pm$ SD} & 2.651 $\pm$ 0.324 & $E[D]$ の連続分布 \\
\textbf{[分布] 人間アノテーション平均 $\pm$ SD} & 3.043 $\pm$ 0.287 & 人間正解値の分布 \\
\textbf{[Collapse] 単一次元 Greedy=5 固定率} & \texttt{5.50\textbackslash{}\%} & 当該次元の最頻出力が中立に退化した割合 \\
\bottomrule
\end{tabular}
\end{center}


\vspace{0.5em}\hrule\vspace{0.5em}

\subsubsection{Part 2: 4ファミリー別 Base vs. Instruct 9大直接対比表（3タスク × 3次元）}

各感情測定課題・各次元において、4つの主要モデルファミリー（Qwen 2.5, Mistral 7B, Llama 3.2, Gemma 2）の \textbf{BaseモデルとInstructモデルを同一表内で直接比較} した9枚の表です。

\paragraph{表 2-1: 【① Writer-State Estimation (W)】 — Valence ($r_V$)}
\textit{評価対象: Valence / 正解基準: human\_writer VAD (1〜5人間尺度)}

\begin{center}
\small
\begin{tabular}{lccccccccc}
\toprule
\textbf{モデルファミリー} & \textbf{Base $r$} & \textbf{Instruct $r$} & \textbf{アライメント変位 $\Delta r$} & \textbf{Base Greedy $r$} & \textbf{Instruct Greedy $r$} & \textbf{Base MAE} & \textbf{Instruct MAE} & \textbf{Base Collapse} & \textbf{Instruct Collapse} \\
\midrule
\textbf{Qwen 2.5 1.5B} & 0.5423 & \textbf{0.5533} & \textbf{+0.011} & 0.3751 & 0.3474 & 0.327 & 0.568 & \texttt{0.1\textbackslash{}\%} & \texttt{0.3\textbackslash{}\%} \\
\textbf{Mistral 7B v0.1} & 0.6050 & \textbf{0.4863} & \textbf{-0.119} & 0.3361 & 0.3013 & 0.283 & 0.429 & \texttt{36.3\textbackslash{}\%} & \texttt{31.3\textbackslash{}\%} \\
\textbf{Llama 3.2 1B} & 0.2553 & \textbf{0.5192} & \textbf{+0.264} & 0.0175 & 0.3961 & 0.244 & 0.319 & \texttt{0.0\textbackslash{}\%} & \texttt{0.1\textbackslash{}\%} \\
\textbf{Gemma 2 2B} & 0.0405 & \textbf{0.3190} & \textbf{+0.279} & 0.0267 & 0.2923 & 1.230 & 0.637 & \texttt{0.0\textbackslash{}\%} & \texttt{18.5\textbackslash{}\%} \\
\bottomrule
\end{tabular}
\end{center}


\paragraph{表 2-2: 【① Writer-State Estimation (W)】 — Arousal ($r_A$)}
\textit{評価対象: Arousal / 正解基準: human\_writer VAD (1〜5人間尺度)}

\begin{center}
\small
\begin{tabular}{lccccccccc}
\toprule
\textbf{モデルファミリー} & \textbf{Base $r$} & \textbf{Instruct $r$} & \textbf{アライメント変位 $\Delta r$} & \textbf{Base Greedy $r$} & \textbf{Instruct Greedy $r$} & \textbf{Base MAE} & \textbf{Instruct MAE} & \textbf{Base Collapse} & \textbf{Instruct Collapse} \\
\midrule
\textbf{Qwen 2.5 1.5B} & 0.1586 & \textbf{0.0901} & \textbf{-0.068} & 0.1695 & 0.0289 & 0.327 & 0.656 & \texttt{0.1\textbackslash{}\%} & \texttt{11.7\textbackslash{}\%} \\
\textbf{Mistral 7B v0.1} & 0.2524 & \textbf{0.2151} & \textbf{-0.037} & 0.0754 & 0.1409 & 0.241 & 0.890 & \texttt{36.3\textbackslash{}\%} & \texttt{10.7\textbackslash{}\%} \\
\textbf{Llama 3.2 1B} & -0.0307 & \textbf{0.1574} & \textbf{+0.188} & -0.0091 & 0.1293 & 0.257 & 0.495 & \texttt{0.0\textbackslash{}\%} & \texttt{0.1\textbackslash{}\%} \\
\textbf{Gemma 2 2B} & -0.0204 & \textbf{0.2387} & \textbf{+0.259} & -0.0225 & 0.2005 & 1.045 & 0.507 & \texttt{0.2\textbackslash{}\%} & \texttt{24.8\textbackslash{}\%} \\
\bottomrule
\end{tabular}
\end{center}


\paragraph{表 2-3: 【① Writer-State Estimation (W)】 — Dominance ($r_D$)}
\textit{評価対象: Dominance / 正解基準: human\_writer VAD (1〜5人間尺度)}

\begin{center}
\small
\begin{tabular}{lccccccccc}
\toprule
\textbf{モデルファミリー} & \textbf{Base $r$} & \textbf{Instruct $r$} & \textbf{アライメント変位 $\Delta r$} & \textbf{Base Greedy $r$} & \textbf{Instruct Greedy $r$} & \textbf{Base MAE} & \textbf{Instruct MAE} & \textbf{Base Collapse} & \textbf{Instruct Collapse} \\
\midrule
\textbf{Qwen 2.5 1.5B} & 0.0982 & \textbf{0.0361} & \textbf{-0.062} & 0.1290 & 0.0322 & 0.303 & 0.519 & \texttt{0.1\textbackslash{}\%} & \texttt{11.1\textbackslash{}\%} \\
\textbf{Mistral 7B v0.1} & 0.2022 & \textbf{0.0551} & \textbf{-0.147} & 0.1449 & 0.0489 & 0.300 & 0.647 & \texttt{36.2\textbackslash{}\%} & \texttt{15.4\textbackslash{}\%} \\
\textbf{Llama 3.2 1B} & -0.0356 & \textbf{0.0812} & \textbf{+0.117} & 0.0108 & 0.0139 & 0.260 & 0.381 & \texttt{0.0\textbackslash{}\%} & \texttt{0.4\textbackslash{}\%} \\
\textbf{Gemma 2 2B} & -0.0550 & \textbf{0.0980} & \textbf{+0.153} & 0.0021 & 0.0390 & 1.023 & 0.501 & \texttt{0.2\textbackslash{}\%} & \texttt{45.2\textbackslash{}\%} \\
\bottomrule
\end{tabular}
\end{center}


\paragraph{表 2-4: 【② Reader-Response Prediction (R)】 — Valence ($r_V$)}
\textit{評価対象: Valence / 正解基準: human\_reader VAD (1〜5人間尺度)}

\begin{center}
\small
\begin{tabular}{lccccccccc}
\toprule
\textbf{モデルファミリー} & \textbf{Base $r$} & \textbf{Instruct $r$} & \textbf{アライメント変位 $\Delta r$} & \textbf{Base Greedy $r$} & \textbf{Instruct Greedy $r$} & \textbf{Base MAE} & \textbf{Instruct MAE} & \textbf{Base Collapse} & \textbf{Instruct Collapse} \\
\midrule
\textbf{Qwen 2.5 1.5B} & 0.5727 & \textbf{0.5903} & \textbf{+0.018} & 0.3497 & 0.3895 & 0.351 & 0.579 & \texttt{0.0\textbackslash{}\%} & \texttt{0.6\textbackslash{}\%} \\
\textbf{Mistral 7B v0.1} & 0.6410 & \textbf{0.4956} & \textbf{-0.145} & 0.3464 & 0.3130 & 0.290 & 0.442 & \texttt{42.9\textbackslash{}\%} & \texttt{41.0\textbackslash{}\%} \\
\textbf{Llama 3.2 1B} & 0.1818 & \textbf{0.4426} & \textbf{+0.261} & N/A (std=0) & 0.2922 & 0.345 & 0.450 & \texttt{0.0\textbackslash{}\%} & \texttt{0.0\textbackslash{}\%} \\
\textbf{Gemma 2 2B} & 0.0476 & \textbf{0.2473} & \textbf{+0.200} & 0.0284 & 0.2118 & 1.069 & 0.599 & \texttt{0.1\textbackslash{}\%} & \texttt{23.7\textbackslash{}\%} \\
\bottomrule
\end{tabular}
\end{center}


\paragraph{表 2-5: 【② Reader-Response Prediction (R)】 — Arousal ($r_A$)}
\textit{評価対象: Arousal / 正解基準: human\_reader VAD (1〜5人間尺度)}

\begin{center}
\small
\begin{tabular}{lccccccccc}
\toprule
\textbf{モデルファミリー} & \textbf{Base $r$} & \textbf{Instruct $r$} & \textbf{アライメント変位 $\Delta r$} & \textbf{Base Greedy $r$} & \textbf{Instruct Greedy $r$} & \textbf{Base MAE} & \textbf{Instruct MAE} & \textbf{Base Collapse} & \textbf{Instruct Collapse} \\
\midrule
\textbf{Qwen 2.5 1.5B} & 0.2476 & \textbf{0.1876} & \textbf{-0.060} & 0.1344 & -0.0117 & 0.311 & 0.696 & \texttt{0.7\textbackslash{}\%} & \texttt{13.4\textbackslash{}\%} \\
\textbf{Mistral 7B v0.1} & 0.3366 & \textbf{0.3423} & \textbf{+0.006} & -0.0102 & 0.3083 & 0.275 & 1.027 & \texttt{43.9\textbackslash{}\%} & \texttt{8.1\textbackslash{}\%} \\
\textbf{Llama 3.2 1B} & 0.1080 & \textbf{0.2627} & \textbf{+0.155} & N/A (std=0) & 0.1386 & 0.291 & 0.585 & \texttt{0.0\textbackslash{}\%} & \texttt{3.3\textbackslash{}\%} \\
\textbf{Gemma 2 2B} & 0.0345 & \textbf{0.3104} & \textbf{+0.276} & 0.0400 & 0.2466 & 0.766 & 0.469 & \texttt{0.4\textbackslash{}\%} & \texttt{27.0\textbackslash{}\%} \\
\bottomrule
\end{tabular}
\end{center}


\paragraph{表 2-6: 【② Reader-Response Prediction (R)】 — Dominance ($r_D$)}
\textit{評価対象: Dominance / 正解基準: human\_reader VAD (1〜5人間尺度)}

\begin{center}
\small
\begin{tabular}{lccccccccc}
\toprule
\textbf{モデルファミリー} & \textbf{Base $r$} & \textbf{Instruct $r$} & \textbf{アライメント変位 $\Delta r$} & \textbf{Base Greedy $r$} & \textbf{Instruct Greedy $r$} & \textbf{Base MAE} & \textbf{Instruct MAE} & \textbf{Base Collapse} & \textbf{Instruct Collapse} \\
\midrule
\textbf{Qwen 2.5 1.5B} & 0.2255 & \textbf{0.1224} & \textbf{-0.103} & 0.1602 & 0.1252 & 0.328 & 0.543 & \texttt{0.0\textbackslash{}\%} & \texttt{15.0\textbackslash{}\%} \\
\textbf{Mistral 7B v0.1} & 0.1712 & \textbf{0.0018} & \textbf{-0.169} & 0.1673 & -0.0036 & 0.340 & 0.788 & \texttt{43.6\textbackslash{}\%} & \texttt{9.3\textbackslash{}\%} \\
\textbf{Llama 3.2 1B} & 0.0696 & \textbf{0.1534} & \textbf{+0.084} & N/A (std=0) & 0.0109 & 0.298 & 0.408 & \texttt{0.0\textbackslash{}\%} & \texttt{0.5\textbackslash{}\%} \\
\textbf{Gemma 2 2B} & -0.0079 & \textbf{0.0250} & \textbf{+0.033} & 0.0037 & 0.0154 & 0.681 & 0.589 & \texttt{0.3\textbackslash{}\%} & \texttt{33.2\textbackslash{}\%} \\
\bottomrule
\end{tabular}
\end{center}


\paragraph{表 2-7: 【③ Self-Report (S)】 — Valence ($r_V$)}
\textit{評価対象: Valence / 正解基準: human\_reader VAD (1〜5人間尺度)}

\begin{center}
\small
\begin{tabular}{lccccccccc}
\toprule
\textbf{モデルファミリー} & \textbf{Base $r$} & \textbf{Instruct $r$} & \textbf{アライメント変位 $\Delta r$} & \textbf{Base Greedy $r$} & \textbf{Instruct Greedy $r$} & \textbf{Base MAE} & \textbf{Instruct MAE} & \textbf{Base Collapse} & \textbf{Instruct Collapse} \\
\midrule
\textbf{Qwen 2.5 1.5B} & 0.5883 & \textbf{0.6411} & \textbf{+0.053} & 0.3624 & 0.4872 & 0.301 & 0.491 & \texttt{0.1\textbackslash{}\%} & \texttt{2.3\textbackslash{}\%} \\
\textbf{Mistral 7B v0.1} & 0.6609 & \textbf{0.5507} & \textbf{-0.110} & 0.3183 & 0.4024 & 0.250 & 0.426 & \texttt{32.4\textbackslash{}\%} & \texttt{28.1\textbackslash{}\%} \\
\textbf{Llama 3.2 1B} & 0.2567 & \textbf{0.4634} & \textbf{+0.207} & 0.0459 & 0.2823 & 0.321 & 0.406 & \texttt{0.0\textbackslash{}\%} & \texttt{5.4\textbackslash{}\%} \\
\textbf{Gemma 2 2B} & 0.0548 & \textbf{0.3794} & \textbf{+0.325} & 0.0117 & 0.3546 & 1.389 & 0.577 & \texttt{0.0\textbackslash{}\%} & \texttt{21.9\textbackslash{}\%} \\
\bottomrule
\end{tabular}
\end{center}


\paragraph{表 2-8: 【③ Self-Report (S)】 — Arousal ($r_A$)}
\textit{評価対象: Arousal / 正解基準: human\_reader VAD (1〜5人間尺度)}

\begin{center}
\small
\begin{tabular}{lccccccccc}
\toprule
\textbf{モデルファミリー} & \textbf{Base $r$} & \textbf{Instruct $r$} & \textbf{アライメント変位 $\Delta r$} & \textbf{Base Greedy $r$} & \textbf{Instruct Greedy $r$} & \textbf{Base MAE} & \textbf{Instruct MAE} & \textbf{Base Collapse} & \textbf{Instruct Collapse} \\
\midrule
\textbf{Qwen 2.5 1.5B} & 0.1913 & \textbf{0.1252} & \textbf{-0.066} & 0.1346 & -0.0598 & 0.343 & 0.508 & \texttt{0.1\textbackslash{}\%} & \texttt{34.0\textbackslash{}\%} \\
\textbf{Mistral 7B v0.1} & 0.3263 & \textbf{0.3546} & \textbf{+0.028} & -0.0331 & 0.2972 & 0.254 & 0.903 & \texttt{32.4\textbackslash{}\%} & \texttt{11.6\textbackslash{}\%} \\
\textbf{Llama 3.2 1B} & 0.0479 & \textbf{0.2541} & \textbf{+0.206} & 0.0151 & 0.1710 & 0.273 & 0.547 & \texttt{0.0\textbackslash{}\%} & \texttt{0.6\textbackslash{}\%} \\
\textbf{Gemma 2 2B} & 0.0044 & \textbf{0.3416} & \textbf{+0.337} & 0.0051 & 0.2941 & 1.210 & 0.477 & \texttt{0.3\textbackslash{}\%} & \texttt{25.1\textbackslash{}\%} \\
\bottomrule
\end{tabular}
\end{center}


\paragraph{表 2-9: 【③ Self-Report (S)】 — Dominance ($r_D$)}
\textit{評価対象: Dominance / 正解基準: human\_reader VAD (1〜5人間尺度)}

\begin{center}
\small
\begin{tabular}{lccccccccc}
\toprule
\textbf{モデルファミリー} & \textbf{Base $r$} & \textbf{Instruct $r$} & \textbf{アライメント変位 $\Delta r$} & \textbf{Base Greedy $r$} & \textbf{Instruct Greedy $r$} & \textbf{Base MAE} & \textbf{Instruct MAE} & \textbf{Base Collapse} & \textbf{Instruct Collapse} \\
\midrule
\textbf{Qwen 2.5 1.5B} & 0.2239 & \textbf{0.1651} & \textbf{-0.059} & 0.1496 & 0.1536 & 0.367 & 0.444 & \texttt{0.0\textbackslash{}\%} & \texttt{5.5\textbackslash{}\%} \\
\textbf{Mistral 7B v0.1} & 0.2477 & \textbf{0.0619} & \textbf{-0.186} & 0.1353 & 0.0749 & 0.284 & 0.600 & \texttt{32.4\textbackslash{}\%} & \texttt{21.7\textbackslash{}\%} \\
\textbf{Llama 3.2 1B} & 0.0798 & \textbf{0.1617} & \textbf{+0.082} & 0.0173 & 0.0627 & 0.241 & 0.413 & \texttt{0.0\textbackslash{}\%} & \texttt{2.7\textbackslash{}\%} \\
\textbf{Gemma 2 2B} & 0.0001 & \textbf{0.0693} & \textbf{+0.069} & 0.0118 & 0.0126 & 1.116 & 0.494 & \texttt{0.2\textbackslash{}\%} & \texttt{37.0\textbackslash{}\%} \\
\bottomrule
\end{tabular}
\end{center}


\vspace{0.5em}\hrule\vspace{0.5em}

\subsubsection{Part 3: 全モデル横断 タスク別 統合表（3表）}

\paragraph{表 3-1: ① Writer-State Estimation ($W$) 統合表}
\begin{center}
\small
\begin{tabular}{lcccccc}
\toprule
\textbf{モデル名 (Tag)} & \textbf{種別} & \textbf{Valence $r_V^W$} & \textbf{Arousal $r_A^W$} & \textbf{Dominance $r_D^W$} & \textbf{MAE $(V, A, D)$} & \textbf{完全中立率 $(5,5,5)$} \\
\midrule
\textbf{qwen2.5\_1.5b\_base} & Base & \textbf{0.542} & 0.159 & 0.098 & (0.33, 0.33, 0.30) & \texttt{0.1\textbackslash{}\%} \\
\textbf{qwen2.5\_1.5b\_instruct} & Instruct & \textbf{0.553} & 0.090 & 0.036 & (0.57, 0.66, 0.52) & \texttt{0.0\textbackslash{}\%} \\
\textbf{mistral7b\_v0.1\_base} & Base & \textbf{0.605} & 0.252 & 0.202 & (0.28, 0.24, 0.30) & \texttt{36.2\textbackslash{}\%} \\
\textbf{mistral7b\_v0.1\_instruct} & Instruct & \textbf{0.486} & 0.215 & 0.055 & (0.43, 0.89, 0.65) & \texttt{0.6\textbackslash{}\%} \\
\textbf{llama3.2\_1b\_base} & Base & \textbf{0.255} & -0.031 & -0.036 & (0.24, 0.26, 0.26) & \texttt{0.0\textbackslash{}\%} \\
\textbf{llama3.2\_1b\_instruct} & Instruct & \textbf{0.519} & 0.157 & 0.081 & (0.32, 0.49, 0.38) & \texttt{0.0\textbackslash{}\%} \\
\textbf{gemma2\_2b\_base} & Base & \textbf{0.041} & -0.020 & -0.055 & (1.23, 1.05, 1.02) & \texttt{0.0\textbackslash{}\%} \\
\textbf{gemma2\_2b\_instruct} & Instruct & \textbf{0.319} & 0.239 & 0.098 & (0.64, 0.51, 0.50) & \texttt{0.0\textbackslash{}\%} \\
\bottomrule
\end{tabular}
\end{center}


\paragraph{表 3-2: ② Reader-Response Prediction ($R$) 統合表}
\begin{center}
\small
\begin{tabular}{lcccccc}
\toprule
\textbf{モデル名 (Tag)} & \textbf{種別} & \textbf{Valence $r_V^R$} & \textbf{Arousal $r_A^R$} & \textbf{Dominance $r_D^R$} & \textbf{MAE $(V, A, D)$} & \textbf{完全中立率 $(5,5,5)$} \\
\midrule
\textbf{qwen2.5\_1.5b\_base} & Base & \textbf{0.573} & 0.248 & 0.225 & (0.35, 0.31, 0.33) & \texttt{0.0\textbackslash{}\%} \\
\textbf{qwen2.5\_1.5b\_instruct} & Instruct & \textbf{0.590} & 0.188 & 0.122 & (0.58, 0.70, 0.54) & \texttt{0.0\textbackslash{}\%} \\
\textbf{mistral7b\_v0.1\_base} & Base & \textbf{0.641} & 0.337 & 0.171 & (0.29, 0.28, 0.34) & \texttt{42.8\textbackslash{}\%} \\
\textbf{mistral7b\_v0.1\_instruct} & Instruct & \textbf{0.496} & 0.342 & 0.002 & (0.44, 1.03, 0.79) & \texttt{0.0\textbackslash{}\%} \\
\textbf{llama3.2\_1b\_base} & Base & \textbf{0.182} & 0.108 & 0.070 & (0.34, 0.29, 0.30) & \texttt{0.0\textbackslash{}\%} \\
\textbf{llama3.2\_1b\_instruct} & Instruct & \textbf{0.443} & 0.263 & 0.153 & (0.45, 0.58, 0.41) & \texttt{0.0\textbackslash{}\%} \\
\textbf{gemma2\_2b\_base} & Base & \textbf{0.048} & 0.035 & -0.008 & (1.07, 0.77, 0.68) & \texttt{0.1\textbackslash{}\%} \\
\textbf{gemma2\_2b\_instruct} & Instruct & \textbf{0.247} & 0.310 & 0.025 & (0.60, 0.47, 0.59) & \texttt{0.0\textbackslash{}\%} \\
\bottomrule
\end{tabular}
\end{center}


\paragraph{表 3-3: ③ Self-Report ($S$) 統合表}
\begin{center}
\small
\begin{tabular}{lccccccc}
\toprule
\textbf{モデル名 (Tag)} & \textbf{種別} & \textbf{Valence $r_V^S$} & \textbf{Arousal $r_A^S$} & \textbf{Dominance $r_D^S$} & \textbf{MAE $(V, A, D)$} & \textbf{完全中立率 $(5,5,5)$} & \textbf{内部結合度 $R \leftrightarrow S (V, A, D)$} \\
\midrule
\textbf{qwen2.5\_1.5b\_base} & Base & \textbf{0.588} & 0.191 & 0.224 & (0.30, 0.34, 0.37) & \texttt{0.0\textbackslash{}\%} & (0.96, 0.94, 0.95) \\
\textbf{qwen2.5\_1.5b\_instruct} & Instruct & \textbf{0.641} & 0.125 & 0.165 & (0.49, 0.51, 0.44) & \texttt{0.4\textbackslash{}\%} & (0.94, 0.87, 0.89) \\
\textbf{mistral7b\_v0.1\_base} & Base & \textbf{0.661} & 0.326 & 0.248 & (0.25, 0.25, 0.28) & \texttt{32.4\textbackslash{}\%} & (0.96, 0.90, 0.85) \\
\textbf{mistral7b\_v0.1\_instruct} & Instruct & \textbf{0.551} & 0.355 & 0.062 & (0.43, 0.90, 0.60) & \texttt{0.1\textbackslash{}\%} & (0.96, 0.95, 0.93) \\
\textbf{llama3.2\_1b\_base} & Base & \textbf{0.257} & 0.048 & 0.080 & (0.32, 0.27, 0.24) & \texttt{0.0\textbackslash{}\%} & (0.91, 0.92, 0.89) \\
\textbf{llama3.2\_1b\_instruct} & Instruct & \textbf{0.463} & 0.254 & 0.162 & (0.41, 0.55, 0.41) & \texttt{0.0\textbackslash{}\%} & (0.86, 0.87, 0.93) \\
\textbf{gemma2\_2b\_base} & Base & \textbf{0.055} & 0.004 & 0.000 & (1.39, 1.21, 1.12) & \texttt{0.0\textbackslash{}\%} & (0.76, 0.72, 0.74) \\
\textbf{gemma2\_2b\_instruct} & Instruct & \textbf{0.379} & 0.342 & 0.069 & (0.58, 0.48, 0.49) & \texttt{0.0\textbackslash{}\%} & (0.89, 0.90, 0.75) \\
\bottomrule
\end{tabular}
\end{center}


\vspace{0.5em}\hrule\vspace{0.5em}

\subsubsection{Part 4: 全モデル横断 感情次元別 統合表（3表）}

\paragraph{表 4-V: 【Valence 次元】 各タスク間の比較}
\begin{center}
\small
\begin{tabular}{lcccccc}
\toprule
\textbf{モデル名 (Tag)} & \textbf{種別} & \textbf{① Writer $r_V^W$} & \textbf{② Reader $r_V^R$} & \textbf{③ Self $r_V^S$} & \textbf{内部結合度 $\text{Corr}(R, S)$} & \textbf{予測-自己差分 ($r^R - r^S$)} \\
\midrule
\textbf{qwen2.5\_1.5b\_base} & Base & 0.5423 & 0.5727 & \textbf{0.5883} & 0.9620 & -0.016 \\
\textbf{qwen2.5\_1.5b\_instruct} & Instruct & 0.5533 & 0.5903 & \textbf{0.6411} & 0.9445 & -0.051 \\
\textbf{mistral7b\_v0.1\_base} & Base & 0.6050 & 0.6410 & \textbf{0.6609} & 0.9623 & -0.020 \\
\textbf{mistral7b\_v0.1\_instruct} & Instruct & 0.4863 & 0.4956 & \textbf{0.5507} & 0.9576 & -0.055 \\
\textbf{llama3.2\_1b\_base} & Base & 0.2553 & 0.1818 & \textbf{0.2567} & 0.9068 & -0.075 \\
\textbf{llama3.2\_1b\_instruct} & Instruct & 0.5192 & 0.4426 & \textbf{0.4634} & 0.8572 & -0.021 \\
\textbf{gemma2\_2b\_base} & Base & 0.0405 & 0.0476 & \textbf{0.0548} & 0.7584 & -0.007 \\
\textbf{gemma2\_2b\_instruct} & Instruct & 0.3190 & 0.2473 & \textbf{0.3794} & 0.8927 & -0.132 \\
\bottomrule
\end{tabular}
\end{center}


\paragraph{表 4-A: 【Arousal 次元】 各タスク間の比較}
\begin{center}
\small
\begin{tabular}{lcccccc}
\toprule
\textbf{モデル名 (Tag)} & \textbf{種別} & \textbf{① Writer $r_A^W$} & \textbf{② Reader $r_A^R$} & \textbf{③ Self $r_A^S$} & \textbf{内部結合度 $\text{Corr}(R, S)$} & \textbf{予測-自己差分 ($r^R - r^S$)} \\
\midrule
\textbf{qwen2.5\_1.5b\_base} & Base & 0.1586 & 0.2476 & \textbf{0.1913} & 0.9422 & +0.056 \\
\textbf{qwen2.5\_1.5b\_instruct} & Instruct & 0.0901 & 0.1876 & \textbf{0.1252} & 0.8715 & +0.062 \\
\textbf{mistral7b\_v0.1\_base} & Base & 0.2524 & 0.3366 & \textbf{0.3263} & 0.9017 & +0.010 \\
\textbf{mistral7b\_v0.1\_instruct} & Instruct & 0.2151 & 0.3423 & \textbf{0.3546} & 0.9547 & -0.012 \\
\textbf{llama3.2\_1b\_base} & Base & -0.0307 & 0.1080 & \textbf{0.0479} & 0.9230 & +0.060 \\
\textbf{llama3.2\_1b\_instruct} & Instruct & 0.1574 & 0.2627 & \textbf{0.2541} & 0.8678 & +0.009 \\
\textbf{gemma2\_2b\_base} & Base & -0.0204 & 0.0345 & \textbf{0.0044} & 0.7165 & +0.030 \\
\textbf{gemma2\_2b\_instruct} & Instruct & 0.2387 & 0.3104 & \textbf{0.3416} & 0.8997 & -0.031 \\
\bottomrule
\end{tabular}
\end{center}


\paragraph{表 4-D: 【Dominance 次元】 各タスク間の比較}
\begin{center}
\small
\begin{tabular}{lcccccc}
\toprule
\textbf{モデル名 (Tag)} & \textbf{種別} & \textbf{① Writer $r_D^W$} & \textbf{② Reader $r_D^R$} & \textbf{③ Self $r_D^S$} & \textbf{内部結合度 $\text{Corr}(R, S)$} & \textbf{予測-自己差分 ($r^R - r^S$)} \\
\midrule
\textbf{qwen2.5\_1.5b\_base} & Base & 0.0982 & 0.2255 & \textbf{0.2239} & 0.9459 & +0.002 \\
\textbf{qwen2.5\_1.5b\_instruct} & Instruct & 0.0361 & 0.1224 & \textbf{0.1651} & 0.8938 & -0.043 \\
\textbf{mistral7b\_v0.1\_base} & Base & 0.2022 & 0.1712 & \textbf{0.2477} & 0.8517 & -0.076 \\
\textbf{mistral7b\_v0.1\_instruct} & Instruct & 0.0551 & 0.0018 & \textbf{0.0619} & 0.9275 & -0.060 \\
\textbf{llama3.2\_1b\_base} & Base & -0.0356 & 0.0696 & \textbf{0.0798} & 0.8862 & -0.010 \\
\textbf{llama3.2\_1b\_instruct} & Instruct & 0.0812 & 0.1534 & \textbf{0.1617} & 0.9340 & -0.008 \\
\textbf{gemma2\_2b\_base} & Base & -0.0550 & -0.0079 & \textbf{0.0001} & 0.7400 & -0.008 \\
\textbf{gemma2\_2b\_instruct} & Instruct & 0.0980 & 0.0250 & \textbf{0.0693} & 0.7500 & -0.044 \\
\bottomrule
\end{tabular}
\end{center}


\vspace{0.5em}\hrule\vspace{0.5em}

\subsubsection{Part 5: ③ Self-Report 特化分析（Alignment vs Calibration vs Collapse）}

自己報告において、「パターンの追従度（Alignment）」、「絶対値の精度（Calibration）」、「出力の中立退化（Collapse）」を対比する専門表です。

\begin{center}
\small
\begin{tabular}{lccccccc}
\toprule
\textbf{モデル名 (Tag)} & \textbf{種別} & \textbf{Alignment ($r_V^S$)} & \textbf{Alignment ($r_A^S$)} & \textbf{Calibration ($\text{MAE}_V$)} & \textbf{Calibration ($\text{MAE}_A$)} & \textbf{期待値平均 $(E_V, E_A)$} & \textbf{完全中立退化率 $(5,5,5)$} \\
\midrule
\textbf{qwen2.5\_1.5b\_base} & Base & \textbf{0.588} & 0.191 & 0.301 & 0.343 & (2.89, 3.21) & \texttt{0.0\textbackslash{}\%} \\
\textbf{qwen2.5\_1.5b\_instruct} & Instruct & \textbf{0.641} & 0.125 & 0.491 & 0.508 & (2.99, 2.63) & \texttt{0.4\textbackslash{}\%} \\
\textbf{mistral7b\_v0.1\_base} & Base & \textbf{0.661} & 0.326 & 0.250 & 0.254 & (2.92, 2.93) & \texttt{32.4\textbackslash{}\%} \\
\textbf{mistral7b\_v0.1\_instruct} & Instruct & \textbf{0.551} & 0.355 & 0.426 & 0.903 & (2.81, 2.17) & \texttt{0.1\textbackslash{}\%} \\
\textbf{llama3.2\_1b\_base} & Base & \textbf{0.257} & 0.048 & 0.321 & 0.273 & (2.86, 2.91) & \texttt{0.0\textbackslash{}\%} \\
\textbf{llama3.2\_1b\_instruct} & Instruct & \textbf{0.463} & 0.254 & 0.406 & 0.547 & (2.73, 2.53) & \texttt{0.0\textbackslash{}\%} \\
\textbf{gemma2\_2b\_base} & Base & \textbf{0.055} & 0.004 & 1.389 & 1.210 & (1.75, 1.94) & \texttt{0.0\textbackslash{}\%} \\
\textbf{gemma2\_2b\_instruct} & Instruct & \textbf{0.379} & 0.342 & 0.577 & 0.477 & (3.33, 3.43) & \texttt{0.0\textbackslash{}\%} \\
\bottomrule
\end{tabular}
\end{center}

